
\documentclass{article}
\usepackage{colortbl}
\usepackage{makecell}
\usepackage{multirow}
\usepackage{supertabular}

\begin{document}

\newcounter{utterance}

\centering \large Interaction Transcript for game `clean\_up', experiment `0\_easy\_5obj\_ga', episode 2 with glm{-}5.2.
\vspace{24pt}

{ \footnotesize  \setcounter{utterance}{1}
\setlength{\tabcolsep}{0pt}
\begin{supertabular}{c@{$\;$}|p{.15\linewidth}@{}p{.15\linewidth}p{.15\linewidth}p{.15\linewidth}p{.15\linewidth}p{.15\linewidth}}
    \# & $\;$A & \multicolumn{4}{c}{Game Master} & $\;\:$B\\
    \hline

    \theutterance \stepcounter{utterance}  
    & & \multicolumn{4}{p{0.6\linewidth}}{
        \cellcolor[rgb]{0.9,0.9,0.9}{
            \makecell[{{p{\linewidth}}}]{
                \texttt{\tiny{[P1$\langle$GM]}}
                \texttt{Is mise do mháistir cluiche, agus tá tú ag imirt cluiche comhoibríoch leis an ngreille seo a leanas mar chlár cluiche:} \\
\\ 
\texttt{```} \\
\texttt{ 1234567} \\
\texttt{╔═══════╗} \\
\texttt{║◌◌◌◌◌◌◌║ 1} \\
\texttt{║◌◌◌◌◌◌◌║ 2} \\
\texttt{║◌TW◌◌┌─╢ 3} \\
\texttt{╟──┐◌◌│H║ 4} \\
\texttt{║C◌├──┼─╢ 5} \\
\texttt{║◌◌│I◌│◌║ 6} \\
\texttt{║◌◌│◌◌│◌║ 7} \\
\texttt{╚══╧══╧═╝} \\
\\ 
\texttt{```} \\
\\ 
\texttt{* Taispeánann an imeall uachtarach comhordanáidí x ag méadú i dtreo na láimhe deise, agus an t{-}imeall ar dheis comhordanáidí y ag méadú síos.} \\
\texttt{* Cuirtear na rudaí seo a leanas go randamach ar do ghreille: 'W', 'I', 'T', 'C', 'H'.} \\
\\ 
\\ 
\texttt{Feiceann an t{-}imreoir eile leagan éagsúil den chlár cluiche, áit a gcuirtear na rudaí in áiteanna randamacha éagsúla. Ní féidir leat clár an imreora eile a fheiceáil, agus ní féidir leo do cheannsa a fheiceáil.} \\
\\ 
\\ 
\texttt{**Sprioc:**} \\
\\ 
\texttt{Caithfidh an dá imreoir na rudaí a bhogadh ar a gcúlra féin ionas go gcríochnóidh rudaí comhionanna ag na comhordanáidí céanna. Caithfidh tú cumarsáid a dhéanamh leis an imreoir eile chun cumraíocht sprice choiteann a aontú.} \\
\\ 
\\ 
\texttt{**Rialacha:**} \\
\\ 
\texttt{* I ngach babhta, is féidir leat ceann amháin go díreach den dá ordú seo a leanas a sheoladh:} \\
\texttt{ 1. `ABAIR: <TEACHTAIREACHT>`: chun teachtaireacht (gach rud suas go dtí an chéad bhriseadh líne eile) a sheoladh chuig an imreoir eile. Cuirfidh mé ar aghaidh í chuig do pháirtí.} \\
\texttt{2. `BOG: <RUD>, (<X>, <Y>)`: chun rud a bhogadh go suíomh nua, áit a bhfuil `<X>` mar an colún agus `<Y>` mar an ró. Cuirfidh mé in iúl duit an raibh do bhogadh rathúil nó nach raibh.} \\
\texttt{* Mura gcloíonn tú leis an bhformáid, nó má sheolann tú roinnt orduithe ag an am céanna, caithfidh mé pionós a ghearradh ort.} \\
\texttt{* Má charnaíonn an dá imreoir níos mó ná 12 pionóis, caillfidh sibh beirt an cluiche.} \\
\texttt{* Tá sé ríthábhachtach go ndéanann tú cumarsáid leis an imreoir eile maidir le do staid sprice! Is é an t{-}aon bhealach *amháin* ar féidir leat do straitéis a chur ar aghaidh chuig an imreoir eile ná an t{-}ordú `ABAIR: <TEACHTAIREACHT>` a úsáid!} \\
\\ 
\\ 
\texttt{**Rudaí a Bhogadh**} \\
\\ 
\texttt{* Ní féidir leat rudaí a bhogadh ach chuig cealla laistigh de theorainneacha na greille. Caithfidh an chill sprice a bheith folamh, is é sin, ní mór di ach an tsiombail '◌' a bheith inti.} \\
\texttt{* Má dhéanann tú iarracht rud a bhogadh go spota nach bhfuil folamh, nó má dhéanann tú iarracht é a bhogadh lasmuigh den ghreille, caithfidh mé pionós a ghearradh ort. Gheobhaidh tú triail eile.} \\
\texttt{* Sula ndéanann tú bogadh, seiceáil faoi dhó go bhfuil an spota sprice folamh, agus nach bhfuil aon litir, fráma ná líne ann!} \\
\\ 
\\ 
\texttt{**Deireadh an Chluiche**} \\
\\ 
\texttt{Má cheapann tú gur shroich tú an sprioc maidir le gach rud a ailíniú, is féidir leat iarraidh ar an imreoir eile an cluiche a chríochnú trí `ABAIR: críochnaithe?` a sheoladh. Má iarrann an t{-}imreoir eile ort an cluiche a chríochnú, agus má fhreagraíonn tú le `ABAIR: críochnaithe!`, tiocfaidh deireadh leis an gcluiche.} \\
\\ 
\texttt{Buaíonn an dá imreoir má chríochnaíonn an cluiche laistigh de 20 babhta, áit a shainmhínítear babhta amháin mar bheirt imreoirí ag seoladh ordú bailí an duine.} \\
\\ 
\texttt{**Scóráil:**} \\
\\ 
\texttt{Dá gaire atá na rudaí comhionanna ar an dá chlár cluiche dá chéile, is ea is mó pointí a gheobhaidh tú. Laghdaíonn pionóis do phointí. An féidir leat an taifead a shárú?} \\
\\ 
\\ 
\texttt{Seol teachtaireacht chuig an imreoir eile chun an cluiche a thosú!} \\
            }
        }
    }
    & \\ \\

    \theutterance \stepcounter{utterance}  
    & \multicolumn{4}{p{0.6\linewidth}}{
        \cellcolor[rgb]{0.9,0.9,0.9}{
            \makecell[{{p{\linewidth}}}]{
                \texttt{\tiny{[P1$\rangle$GM]}}
                \texttt{ABAIR: Dia duit! Tá mé réidh le tús a chur leis an gcluiche. Fulaingím mo chlár agus cuirim na 5 rud (W, I, T, C, H) in ionaid éagsúla uaitse. Chun críche comhoibríoch a shocrú go héasca agus\_pointí a uasmhéadú, molaim an staid shprice seo a leanas duit:} \\
\\ 
\texttt{Sprioc A:} \\
\texttt{* W ag (1, 6)} \\
\texttt{* I ag (2, 7)} \\
\texttt{* T ag (3, 6)} \\
\texttt{* C ag (4, 7)} \\
\texttt{* H ag (5, 6)} \\
\\ 
\texttt{Fan go dtí go ndeir tú " ullamh " sula dtosóidh muid ag bogadh!} \\
            }
        }
    }
    & & \\ \\

    \theutterance \stepcounter{utterance}  
    & & & \multicolumn{2}{p{0.3\linewidth}}{
        \cellcolor[rgb]{0.9,0.9,0.9}{
            \makecell[{{p{\linewidth}}}]{
                \texttt{\tiny{[GM$|$GM]}}
                \texttt{ABAIR: Dia duit! Tá mé réidh le tús a chur leis an gcluiche. Fulaingím mo chlár agus cuirim na 5 rud (W, I, T, C, H) in ionaid éagsúla uaitse. Chun críche comhoibríoch a shocrú go héasca agus\_pointí a uasmhéadú, molaim an staid shprice seo a leanas duit:} \\
\\ 
\texttt{Sprioc A:} \\
\texttt{* W ag (1, 6)} \\
\texttt{* I ag (2, 7)} \\
\texttt{* T ag (3, 6)} \\
\texttt{* C ag (4, 7)} \\
\texttt{* H ag (5, 6)} \\
\\ 
\texttt{Fan go dtí go ndeir tú " ullamh " sula dtosóidh muid ag bogadh!} \\
            }
        }
    }
    & & \\ \\

    \theutterance \stepcounter{utterance}  
    & & & \multicolumn{2}{p{0.3\linewidth}}{
        \cellcolor[rgb]{0.9,0.9,0.9}{
            \makecell[{{p{\linewidth}}}]{
                \texttt{\tiny{[GM$|$GM]}}
                \texttt{Invalid format: tail is not empty: '} \\
\\ 
\texttt{Sprioc A:} \\
\texttt{* W ag (1, 6)} \\
\texttt{* I ag (2, 7)} \\
\texttt{* T ag (3, 6)} \\
\texttt{* C ag (4, 7)} \\
\texttt{* H ag (5, 6)} \\
\\ 
\texttt{Fan go dtí go ndeir tú " ullamh " sula dtosóidh muid ag bogadh!} \\
            }
        }
    }
    & & \\ \\

    \theutterance \stepcounter{utterance}  
    & & \multicolumn{4}{p{0.6\linewidth}}{
        \cellcolor[rgb]{0.9,0.9,0.9}{
            \makecell[{{p{\linewidth}}}]{
                \texttt{\tiny{[P1$\langle$GM]}}
                \texttt{Pionós: Ní cheadaítear aon rud a bheith i do theachtaireacht tar éis an ordaithe!} \\
\\ 
\texttt{Cinntigh nach bhfuil i do fhreagra ach `ABAIR: <TEACHTAIREACHT>` nó `BOG: <RUD>, (<X>, <Y>)`, agus tada eile!} \\
\\ 
\texttt{Tá 1 as 12 pionós carntha agaibh le chéile.} \\
\texttt{Déan iarracht arís, le do thoil!} \\
            }
        }
    }
    & \\ \\

    \theutterance \stepcounter{utterance}  
    & \multicolumn{4}{p{0.6\linewidth}}{
        \cellcolor[rgb]{0.9,0.9,0.9}{
            \makecell[{{p{\linewidth}}}]{
                \texttt{\tiny{[P1$\rangle$GM]}}
                \texttt{ABAIR: Dia duit! Seo iad na comhordanáidí atá agamsa: T ag 3,3, W ag 4,3, H ag 7,4, C ag 1,5, I ag 3,6. Cén chuma ar do chlár? Molaim sprioc a shocrú. Mar shampla: W ag 1,6, I ag 2,7, T ag 3,6, C ag 4,7, H ag 5,6. An bhfuil tú sásta leis sin?} \\
            }
        }
    }
    & & \\ \\

    \theutterance \stepcounter{utterance}  
    & & & \multicolumn{2}{p{0.3\linewidth}}{
        \cellcolor[rgb]{0.9,0.9,0.9}{
            \makecell[{{p{\linewidth}}}]{
                \texttt{\tiny{[GM$|$GM]}}
                \texttt{ABAIR: Dia duit! Seo iad na comhordanáidí atá agamsa: T ag 3,3, W ag 4,3, H ag 7,4, C ag 1,5, I ag 3,6. Cén chuma ar do chlár? Molaim sprioc a shocrú. Mar shampla: W ag 1,6, I ag 2,7, T ag 3,6, C ag 4,7, H ag 5,6. An bhfuil tú sásta leis sin?} \\
            }
        }
    }
    & & \\ \\

    \theutterance \stepcounter{utterance}  
    & & \multicolumn{4}{p{0.6\linewidth}}{
        \cellcolor[rgb]{0.9,0.9,0.9}{
            \makecell[{{p{\linewidth}}}]{
                \texttt{\tiny{[GM$\rangle$P2]}}
                \texttt{Is mise do mháistir cluiche, agus tá tú ag imirt cluiche comhoibríoch leis an ngreille seo a leanas mar chlár cluiche:} \\
\\ 
\texttt{```} \\
\texttt{ 1234567} \\
\texttt{╔═══════╗} \\
\texttt{║◌◌◌◌TH◌║ 1} \\
\texttt{║◌◌◌◌◌◌◌║ 2} \\
\texttt{║◌◌◌◌◌┌─╢ 3} \\
\texttt{╟──┐◌◌│◌║ 4} \\
\texttt{║I◌├──┼─╢ 5} \\
\texttt{║◌W│◌◌│C║ 6} \\
\texttt{║◌◌│◌◌│◌║ 7} \\
\texttt{╚══╧══╧═╝} \\
\\ 
\texttt{```} \\
\\ 
\texttt{* Taispeánann an imeall uachtarach comhordanáidí x ag méadú i dtreo na láimhe deise, agus an t{-}imeall ar dheis comhordanáidí y ag méadú síos.} \\
\texttt{* Cuirtear na rudaí seo a leanas go randamach ar do ghreille: 'W', 'I', 'T', 'C', 'H'.} \\
\\ 
\\ 
\texttt{Feiceann an t{-}imreoir eile leagan éagsúil den chlár cluiche, áit a gcuirtear na rudaí in áiteanna randamacha éagsúla. Ní féidir leat clár an imreora eile a fheiceáil, agus ní féidir leo do cheannsa a fheiceáil.} \\
\\ 
\\ 
\texttt{**Sprioc:**} \\
\\ 
\texttt{Caithfidh an dá imreoir na rudaí a bhogadh ar a gcúlra féin ionas go gcríochnóidh rudaí comhionanna ag na comhordanáidí céanna. Caithfidh tú cumarsáid a dhéanamh leis an imreoir eile chun cumraíocht sprice choiteann a aontú.} \\
\\ 
\\ 
\texttt{**Rialacha:**} \\
\\ 
\texttt{* I ngach babhta, is féidir leat ceann amháin go díreach den dá ordú seo a leanas a sheoladh:} \\
\texttt{ 1. `ABAIR: <TEACHTAIREACHT>`: chun teachtaireacht (gach rud suas go dtí an chéad bhriseadh líne eile) a sheoladh chuig an imreoir eile. Cuirfidh mé ar aghaidh í chuig do pháirtí.} \\
\texttt{2. `BOG: <RUD>, (<X>, <Y>)`: chun rud a bhogadh go suíomh nua, áit a bhfuil `<X>` mar an colún agus `<Y>` mar an ró. Cuirfidh mé in iúl duit an raibh do bhogadh rathúil nó nach raibh.} \\
\texttt{* Mura gcloíonn tú leis an bhformáid, nó má sheolann tú roinnt orduithe ag an am céanna, caithfidh mé pionós a ghearradh ort.} \\
\texttt{* Má charnaíonn an dá imreoir níos mó ná 12 pionóis, caillfidh sibh beirt an cluiche.} \\
\texttt{* Tá sé ríthábhachtach go ndéanann tú cumarsáid leis an imreoir eile maidir le do staid sprice! Is é an t{-}aon bhealach *amháin* ar féidir leat do straitéis a chur ar aghaidh chuig an imreoir eile ná an t{-}ordú `ABAIR: <TEACHTAIREACHT>` a úsáid!} \\
\\ 
\\ 
\texttt{**Rudaí a Bhogadh**} \\
\\ 
\texttt{* Ní féidir leat rudaí a bhogadh ach chuig cealla laistigh de theorainneacha na greille. Caithfidh an chill sprice a bheith folamh, is é sin, ní mór di ach an tsiombail '◌' a bheith inti.} \\
\texttt{* Má dhéanann tú iarracht rud a bhogadh go spota nach bhfuil folamh, nó má dhéanann tú iarracht é a bhogadh lasmuigh den ghreille, caithfidh mé pionós a ghearradh ort. Gheobhaidh tú triail eile.} \\
\texttt{* Sula ndéanann tú bogadh, seiceáil faoi dhó go bhfuil an spota sprice folamh, agus nach bhfuil aon litir, fráma ná líne ann!} \\
\\ 
\\ 
\texttt{**Deireadh an Chluiche**} \\
\\ 
\texttt{Má cheapann tú gur shroich tú an sprioc maidir le gach rud a ailíniú, is féidir leat iarraidh ar an imreoir eile an cluiche a chríochnú trí `ABAIR: críochnaithe?` a sheoladh. Má iarrann an t{-}imreoir eile ort an cluiche a chríochnú, agus má fhreagraíonn tú le `ABAIR: críochnaithe!`, tiocfaidh deireadh leis an gcluiche.} \\
\\ 
\texttt{Buaíonn an dá imreoir má chríochnaíonn an cluiche laistigh de 20 babhta, áit a shainmhínítear babhta amháin mar bheirt imreoirí ag seoladh ordú bailí an duine.} \\
\\ 
\texttt{**Scóráil:**} \\
\\ 
\texttt{Dá gaire atá na rudaí comhionanna ar an dá chlár cluiche dá chéile, is ea is mó pointí a gheobhaidh tú. Laghdaíonn pionóis do phointí. An féidir leat an taifead a shárú?} \\
\\ 
\texttt{Thosaigh an t{-}imreoir eile an cluiche trí an teachtaireacht seo a sheoladh:} \\
\\ 
\texttt{"Dia duit! Seo iad na comhordanáidí atá agamsa: T ag 3,3, W ag 4,3, H ag 7,4, C ag 1,5, I ag 3,6. Cén chuma ar do chlár? Molaim sprioc a shocrú. Mar shampla: W ag 1,6, I ag 2,7, T ag 3,6, C ag 4,7, H ag 5,6. An bhfuil tú sásta leis sin?"} \\
\\ 
\texttt{Cad é do chéad ordú?} \\
            }
        }
    }
    & \\ \\

    \theutterance \stepcounter{utterance}  
    & & & \multicolumn{4}{p{0.6\linewidth}}{
        \cellcolor[rgb]{0.9,0.9,0.9}{
            \makecell[{{p{\linewidth}}}]{
                \texttt{\tiny{[GM$\langle$P2]}}
                \texttt{ABAIR: Dia duit! Seo í mo chlár: W ag (2,6), I ag (2,5), T ag (6,3), C ag (7,6), H ag (6,1). Tá an chuma ar an scéal gur beagán cathrach atá againn le hathrú. Tá do phropósadh go breá ach iarrfaidinn orainn sprioc atá níos cóirí don dá thaobh a phlé. Tá na cealla folmhacha seo agam: (1,1) go (7,7) seachas na spotaí a bhfuil frámaí nó mo chuid rudaí iontu cheana féin. D'fhéadfaimis an sprioc a dhéanamh ag: I ag (2,2), W ag (2,3), T ag (6,6), C ag (7,7), H ag (6,5). Cad é do bharúil?} \\
            }
        }
    }
     \\ \\

    \theutterance \stepcounter{utterance}  
    & & & \multicolumn{2}{p{0.3\linewidth}}{
        \cellcolor[rgb]{0.9,0.9,0.9}{
            \makecell[{{p{\linewidth}}}]{
                \texttt{\tiny{[GM$|$GM]}}
                \texttt{ABAIR: Dia duit! Seo í mo chlár: W ag (2,6), I ag (2,5), T ag (6,3), C ag (7,6), H ag (6,1). Tá an chuma ar an scéal gur beagán cathrach atá againn le hathrú. Tá do phropósadh go breá ach iarrfaidinn orainn sprioc atá níos cóirí don dá thaobh a phlé. Tá na cealla folmhacha seo agam: (1,1) go (7,7) seachas na spotaí a bhfuil frámaí nó mo chuid rudaí iontu cheana féin. D'fhéadfaimis an sprioc a dhéanamh ag: I ag (2,2), W ag (2,3), T ag (6,6), C ag (7,7), H ag (6,5). Cad é do bharúil?} \\
            }
        }
    }
    & & \\ \\

    \theutterance \stepcounter{utterance}  
    & & \multicolumn{4}{p{0.6\linewidth}}{
        \cellcolor[rgb]{0.9,0.9,0.9}{
            \makecell[{{p{\linewidth}}}]{
                \texttt{\tiny{[P1$\langle$GM]}}
                \texttt{Cuireadh do theachtaireacht ar aghaidh chuig an imreoir eile.} \\
\\ 
\\ 
\texttt{Tá sibh ag imirt babhta 1 faoi láthair as uasmhéid 20.} \\
\texttt{Tá 1 as 12 pionós carntha agaibh le chéile.} \\
\texttt{Sheol an t{-}imreoir eile an teachtaireacht seo:} \\
\\ 
\\ 
\texttt{"Dia duit! Seo í mo chlár: W ag (2,6), I ag (2,5), T ag (6,3), C ag (7,6), H ag (6,1). Tá an chuma ar an scéal gur beagán cathrach atá againn le hathrú. Tá do phropósadh go breá ach iarrfaidinn orainn sprioc atá níos cóirí don dá thaobh a phlé. Tá na cealla folmhacha seo agam: (1,1) go (7,7) seachas na spotaí a bhfuil frámaí nó mo chuid rudaí iontu cheana féin. D'fhéadfaimis an sprioc a dhéanamh ag: I ag (2,2), W ag (2,3), T ag (6,6), C ag (7,7), H ag (6,5). Cad é do bharúil?"} \\
\\ 
\\ 
\texttt{Cad é do chéad ordú eile?} \\
            }
        }
    }
    & \\ \\

    \theutterance \stepcounter{utterance}  
    & \multicolumn{4}{p{0.6\linewidth}}{
        \cellcolor[rgb]{0.9,0.9,0.9}{
            \makecell[{{p{\linewidth}}}]{
                \texttt{\tiny{[P1$\rangle$GM]}}
                \texttt{ABAIR: Tá do sprioc an{-}mhaith go hiomlán! Glacaim leis. Tosóidh mé ag bogadh anois. Beidh mise ag bogadh T go (6,6) agus W go (2,3). Abair liom nuair atá do bhogadh déanta agus inis dom céard atá fágaithe agat.} \\
            }
        }
    }
    & & \\ \\

    \theutterance \stepcounter{utterance}  
    & & & \multicolumn{2}{p{0.3\linewidth}}{
        \cellcolor[rgb]{0.9,0.9,0.9}{
            \makecell[{{p{\linewidth}}}]{
                \texttt{\tiny{[GM$|$GM]}}
                \texttt{ABAIR: Tá do sprioc an{-}mhaith go hiomlán! Glacaim leis. Tosóidh mé ag bogadh anois. Beidh mise ag bogadh T go (6,6) agus W go (2,3). Abair liom nuair atá do bhogadh déanta agus inis dom céard atá fágaithe agat.} \\
            }
        }
    }
    & & \\ \\

    \theutterance \stepcounter{utterance}  
    & & \multicolumn{4}{p{0.6\linewidth}}{
        \cellcolor[rgb]{0.9,0.9,0.9}{
            \makecell[{{p{\linewidth}}}]{
                \texttt{\tiny{[GM$\rangle$P2]}}
                \texttt{Cuireadh do theachtaireacht ar aghaidh chuig an imreoir eile.} \\
\\ 
\\ 
\texttt{Tá sibh ag imirt babhta 2 faoi láthair as uasmhéid 20.} \\
\texttt{Tá 1 as 12 pionós carntha agaibh le chéile.} \\
\texttt{Sheol an t{-}imreoir eile an teachtaireacht seo:} \\
\\ 
\\ 
\texttt{"Tá do sprioc an{-}mhaith go hiomlán! Glacaim leis. Tosóidh mé ag bogadh anois. Beidh mise ag bogadh T go (6,6) agus W go (2,3). Abair liom nuair atá do bhogadh déanta agus inis dom céard atá fágaithe agat."} \\
\\ 
\\ 
\texttt{Cad é do chéad ordú eile?} \\
            }
        }
    }
    & \\ \\

    \theutterance \stepcounter{utterance}  
    & & & \multicolumn{4}{p{0.6\linewidth}}{
        \cellcolor[rgb]{0.9,0.9,0.9}{
            \makecell[{{p{\linewidth}}}]{
                \texttt{\tiny{[GM$\langle$P2]}}
                \texttt{ABAIR: Táimid ar aon fhocal! Bainfidh mé triail as I a bhogadh go (2,2) anois. Fágfaidh sin H ag (6,1) agus C ag (7,6) liom le bogadh níos déanaí. Inis dom conas a d'éirigh le do bhogaithe nuair a bheidh siad déanta.} \\
            }
        }
    }
     \\ \\

    \theutterance \stepcounter{utterance}  
    & & & \multicolumn{2}{p{0.3\linewidth}}{
        \cellcolor[rgb]{0.9,0.9,0.9}{
            \makecell[{{p{\linewidth}}}]{
                \texttt{\tiny{[GM$|$GM]}}
                \texttt{ABAIR: Táimid ar aon fhocal! Bainfidh mé triail as I a bhogadh go (2,2) anois. Fágfaidh sin H ag (6,1) agus C ag (7,6) liom le bogadh níos déanaí. Inis dom conas a d'éirigh le do bhogaithe nuair a bheidh siad déanta.} \\
            }
        }
    }
    & & \\ \\

    \theutterance \stepcounter{utterance}  
    & & \multicolumn{4}{p{0.6\linewidth}}{
        \cellcolor[rgb]{0.9,0.9,0.9}{
            \makecell[{{p{\linewidth}}}]{
                \texttt{\tiny{[P1$\langle$GM]}}
                \texttt{Cuireadh do theachtaireacht ar aghaidh chuig an imreoir eile.} \\
\\ 
\\ 
\texttt{Tá sibh ag imirt babhta 2 faoi láthair as uasmhéid 20.} \\
\texttt{Tá 1 as 12 pionós carntha agaibh le chéile.} \\
\texttt{Sheol an t{-}imreoir eile an teachtaireacht seo:} \\
\\ 
\\ 
\texttt{"Táimid ar aon fhocal! Bainfidh mé triail as I a bhogadh go (2,2) anois. Fágfaidh sin H ag (6,1) agus C ag (7,6) liom le bogadh níos déanaí. Inis dom conas a d'éirigh le do bhogaithe nuair a bheidh siad déanta."} \\
\\ 
\\ 
\texttt{Cad é do chéad ordú eile?} \\
            }
        }
    }
    & \\ \\

    \theutterance \stepcounter{utterance}  
    & \multicolumn{4}{p{0.6\linewidth}}{
        \cellcolor[rgb]{0.9,0.9,0.9}{
            \makecell[{{p{\linewidth}}}]{
                \texttt{\tiny{[P1$\rangle$GM]}}
                \texttt{BOG: T, (6, 6)} \\
            }
        }
    }
    & & \\ \\

    \theutterance \stepcounter{utterance}  
    & & & \multicolumn{2}{p{0.3\linewidth}}{
        \cellcolor[rgb]{0.9,0.9,0.9}{
            \makecell[{{p{\linewidth}}}]{
                \texttt{\tiny{[GM$|$GM]}}
                \texttt{BOG: T, (6, 6)} \\
            }
        }
    }
    & & \\ \\

    \theutterance \stepcounter{utterance}  
    & & & \multicolumn{2}{p{0.3\linewidth}}{
        \cellcolor[rgb]{0.9,0.9,0.9}{
            \makecell[{{p{\linewidth}}}]{
                \texttt{\tiny{[GM$|$GM]}}
                \texttt{Pionós: níl (6,6) folamh, ach tá '│' ann.} \\
\\ 
\texttt{Tá 2 as 12 pionós carntha agaibh le chéile.} \\
\texttt{Déan iarracht arís, le do thoil!} \\
            }
        }
    }
    & & \\ \\

    \theutterance \stepcounter{utterance}  
    & & \multicolumn{4}{p{0.6\linewidth}}{
        \cellcolor[rgb]{0.9,0.9,0.9}{
            \makecell[{{p{\linewidth}}}]{
                \texttt{\tiny{[P1$\langle$GM]}}
                \texttt{Pionós: níl (6,6) folamh, ach tá '│' ann.} \\
\\ 
\texttt{Tá 2 as 12 pionós carntha agaibh le chéile.} \\
\texttt{Déan iarracht arís, le do thoil!} \\
            }
        }
    }
    & \\ \\

    \theutterance \stepcounter{utterance}  
    & \multicolumn{4}{p{0.6\linewidth}}{
        \cellcolor[rgb]{0.9,0.9,0.9}{
            \makecell[{{p{\linewidth}}}]{
                \texttt{\tiny{[P1$\rangle$GM]}}
                \texttt{ABAIR: Ó! Bhí fráma ag (6,6) ar mo chlár, ní folamh. Níorbh fhéidir liom T a bhogadh. Tá mé chun iarracht a dhéanamh in ionad éigin eile, b'fhéidir (1,1). Seiceálfaidh mé. Céard é do staid anois?} \\
            }
        }
    }
    & & \\ \\

    \theutterance \stepcounter{utterance}  
    & & & \multicolumn{2}{p{0.3\linewidth}}{
        \cellcolor[rgb]{0.9,0.9,0.9}{
            \makecell[{{p{\linewidth}}}]{
                \texttt{\tiny{[GM$|$GM]}}
                \texttt{ABAIR: Ó! Bhí fráma ag (6,6) ar mo chlár, ní folamh. Níorbh fhéidir liom T a bhogadh. Tá mé chun iarracht a dhéanamh in ionad éigin eile, b'fhéidir (1,1). Seiceálfaidh mé. Céard é do staid anois?} \\
            }
        }
    }
    & & \\ \\

    \theutterance \stepcounter{utterance}  
    & & \multicolumn{4}{p{0.6\linewidth}}{
        \cellcolor[rgb]{0.9,0.9,0.9}{
            \makecell[{{p{\linewidth}}}]{
                \texttt{\tiny{[GM$\rangle$P2]}}
                \texttt{Cuireadh do theachtaireacht ar aghaidh chuig an imreoir eile.} \\
\\ 
\\ 
\texttt{Tá sibh ag imirt babhta 3 faoi láthair as uasmhéid 20.} \\
\texttt{Tá 2 as 12 pionós carntha agaibh le chéile.} \\
\texttt{Sheol an t{-}imreoir eile an teachtaireacht seo:} \\
\\ 
\\ 
\texttt{"Ó! Bhí fráma ag (6,6) ar mo chlár, ní folamh. Níorbh fhéidir liom T a bhogadh. Tá mé chun iarracht a dhéanamh in ionad éigin eile, b'fhéidir (1,1). Seiceálfaidh mé. Céard é do staid anois?"} \\
\\ 
\\ 
\texttt{Cad é do chéad ordú eile?} \\
            }
        }
    }
    & \\ \\

    \theutterance \stepcounter{utterance}  
    & & & \multicolumn{4}{p{0.6\linewidth}}{
        \cellcolor[rgb]{0.9,0.9,0.9}{
            \makecell[{{p{\linewidth}}}]{
                \texttt{\tiny{[GM$\langle$P2]}}
                \texttt{ABAIR: Níor mhaith liom pionós a fháil, mar sin fanfaidh mé go fóill. Níor bhog mé féin rud ar bith fós sa bhabhta seo mar gheall air sin. Tá mo sprioc fós soiléir: W ag (2,3), T ag (6,6), I ag (2,2), C ag (7,7), H ag (6,5). Inis dom conas a d'éirigh leat le T ag (1,1).} \\
            }
        }
    }
     \\ \\

    \theutterance \stepcounter{utterance}  
    & & & \multicolumn{2}{p{0.3\linewidth}}{
        \cellcolor[rgb]{0.9,0.9,0.9}{
            \makecell[{{p{\linewidth}}}]{
                \texttt{\tiny{[GM$|$GM]}}
                \texttt{ABAIR: Níor mhaith liom pionós a fháil, mar sin fanfaidh mé go fóill. Níor bhog mé féin rud ar bith fós sa bhabhta seo mar gheall air sin. Tá mo sprioc fós soiléir: W ag (2,3), T ag (6,6), I ag (2,2), C ag (7,7), H ag (6,5). Inis dom conas a d'éirigh leat le T ag (1,1).} \\
            }
        }
    }
    & & \\ \\

    \theutterance \stepcounter{utterance}  
    & & \multicolumn{4}{p{0.6\linewidth}}{
        \cellcolor[rgb]{0.9,0.9,0.9}{
            \makecell[{{p{\linewidth}}}]{
                \texttt{\tiny{[P1$\langle$GM]}}
                \texttt{Cuireadh do theachtaireacht ar aghaidh chuig an imreoir eile.} \\
\\ 
\\ 
\texttt{Tá sibh ag imirt babhta 3 faoi láthair as uasmhéid 20.} \\
\texttt{Tá 2 as 12 pionós carntha agaibh le chéile.} \\
\texttt{Sheol an t{-}imreoir eile an teachtaireacht seo:} \\
\\ 
\\ 
\texttt{"Níor mhaith liom pionós a fháil, mar sin fanfaidh mé go fóill. Níor bhog mé féin rud ar bith fós sa bhabhta seo mar gheall air sin. Tá mo sprioc fós soiléir: W ag (2,3), T ag (6,6), I ag (2,2), C ag (7,7), H ag (6,5). Inis dom conas a d'éirigh leat le T ag (1,1)."} \\
\\ 
\\ 
\texttt{Cad é do chéad ordú eile?} \\
            }
        }
    }
    & \\ \\

    \theutterance \stepcounter{utterance}  
    & \multicolumn{4}{p{0.6\linewidth}}{
        \cellcolor[rgb]{0.9,0.9,0.9}{
            \makecell[{{p{\linewidth}}}]{
                \texttt{\tiny{[P1$\rangle$GM]}}
                \texttt{BOG: T, (1, 1)} \\
            }
        }
    }
    & & \\ \\

    \theutterance \stepcounter{utterance}  
    & & & \multicolumn{2}{p{0.3\linewidth}}{
        \cellcolor[rgb]{0.9,0.9,0.9}{
            \makecell[{{p{\linewidth}}}]{
                \texttt{\tiny{[GM$|$GM]}}
                \texttt{BOG: T, (1, 1)} \\
            }
        }
    }
    & & \\ \\

    \theutterance \stepcounter{utterance}  
    & & & \multicolumn{2}{p{0.3\linewidth}}{
        \cellcolor[rgb]{0.9,0.9,0.9}{
            \makecell[{{p{\linewidth}}}]{
                \texttt{\tiny{[GM$|$GM]}}
                \texttt{T moved to (1, 1)} \\
            }
        }
    }
    & & \\ \\

    \theutterance \stepcounter{utterance}  
    & & \multicolumn{4}{p{0.6\linewidth}}{
        \cellcolor[rgb]{0.9,0.9,0.9}{
            \makecell[{{p{\linewidth}}}]{
                \texttt{\tiny{[GM$\rangle$P2]}}
                \texttt{Cuireadh do theachtaireacht ar aghaidh chuig an imreoir eile.} \\
\\ 
\\ 
\texttt{Tá sibh ag imirt babhta 4 faoi láthair as uasmhéid 20.} \\
\texttt{Tá 2 as 12 pionós carntha agaibh le chéile.} \\
\texttt{Bhog an t{-}imreoir eile réad ar a ngreille.} \\
\\ 
\\ 
\texttt{Cad é do chéad ordú eile?} \\
            }
        }
    }
    & \\ \\

    \theutterance \stepcounter{utterance}  
    & & & \multicolumn{4}{p{0.6\linewidth}}{
        \cellcolor[rgb]{0.9,0.9,0.9}{
            \makecell[{{p{\linewidth}}}]{
                \texttt{\tiny{[GM$\langle$P2]}}
                \texttt{BOG: I, (2, 2)} \\
            }
        }
    }
     \\ \\

    \theutterance \stepcounter{utterance}  
    & & & \multicolumn{2}{p{0.3\linewidth}}{
        \cellcolor[rgb]{0.9,0.9,0.9}{
            \makecell[{{p{\linewidth}}}]{
                \texttt{\tiny{[GM$|$GM]}}
                \texttt{BOG: I, (2, 2)} \\
            }
        }
    }
    & & \\ \\

    \theutterance \stepcounter{utterance}  
    & & & \multicolumn{2}{p{0.3\linewidth}}{
        \cellcolor[rgb]{0.9,0.9,0.9}{
            \makecell[{{p{\linewidth}}}]{
                \texttt{\tiny{[GM$|$GM]}}
                \texttt{I moved to (2, 2)} \\
            }
        }
    }
    & & \\ \\

    \theutterance \stepcounter{utterance}  
    & & \multicolumn{4}{p{0.6\linewidth}}{
        \cellcolor[rgb]{0.9,0.9,0.9}{
            \makecell[{{p{\linewidth}}}]{
                \texttt{\tiny{[P1$\langle$GM]}}
                \texttt{Bogadh 'T' go (1,1) go rathúil. Seo mar atá do ghreille nuashonraithe:} \\
\texttt{```} \\
\texttt{ 1234567} \\
\texttt{╔═══════╗} \\
\texttt{║T◌◌◌◌◌◌║ 1} \\
\texttt{║◌◌◌◌◌◌◌║ 2} \\
\texttt{║◌◌W◌◌┌─╢ 3} \\
\texttt{╟──┐◌◌│H║ 4} \\
\texttt{║C◌├──┼─╢ 5} \\
\texttt{║◌◌│I◌│◌║ 6} \\
\texttt{║◌◌│◌◌│◌║ 7} \\
\texttt{╚══╧══╧═╝} \\
\texttt{```} \\
\texttt{Tá sibh ag imirt babhta 4 faoi láthair as uasmhéid 20.} \\
\texttt{Tá 2 as 12 pionós carntha agaibh le chéile.} \\
\texttt{Bhog an t{-}imreoir eile réad ar a ngreille.} \\
\\ 
\\ 
\texttt{Cad é do chéad ordú eile?} \\
            }
        }
    }
    & \\ \\

    \theutterance \stepcounter{utterance}  
    & \multicolumn{4}{p{0.6\linewidth}}{
        \cellcolor[rgb]{0.9,0.9,0.9}{
            \makecell[{{p{\linewidth}}}]{
                \texttt{\tiny{[P1$\rangle$GM]}}
                \texttt{ABAIR: D'éirigh liom T a bhogadh go (1,1). Ach tá frámaí ag (6,6) agus (7,7) ar mo chlár, mar sin ní féidir liom spriocanna a gcuid a shroichint ansin. Molaim sprioc nua: I ag (2,6), C ag (3,7), H ag (4,6). Cad é do bharúil?} \\
            }
        }
    }
    & & \\ \\

    \theutterance \stepcounter{utterance}  
    & & & \multicolumn{2}{p{0.3\linewidth}}{
        \cellcolor[rgb]{0.9,0.9,0.9}{
            \makecell[{{p{\linewidth}}}]{
                \texttt{\tiny{[GM$|$GM]}}
                \texttt{ABAIR: D'éirigh liom T a bhogadh go (1,1). Ach tá frámaí ag (6,6) agus (7,7) ar mo chlár, mar sin ní féidir liom spriocanna a gcuid a shroichint ansin. Molaim sprioc nua: I ag (2,6), C ag (3,7), H ag (4,6). Cad é do bharúil?} \\
            }
        }
    }
    & & \\ \\

    \theutterance \stepcounter{utterance}  
    & & \multicolumn{4}{p{0.6\linewidth}}{
        \cellcolor[rgb]{0.9,0.9,0.9}{
            \makecell[{{p{\linewidth}}}]{
                \texttt{\tiny{[GM$\rangle$P2]}}
                \texttt{Bogadh 'I' go (2,2) go rathúil. Seo mar atá do ghreille nuashonraithe:} \\
\texttt{```} \\
\texttt{ 1234567} \\
\texttt{╔═══════╗} \\
\texttt{║◌◌◌◌TH◌║ 1} \\
\texttt{║◌I◌◌◌◌◌║ 2} \\
\texttt{║◌◌◌◌◌┌─╢ 3} \\
\texttt{╟──┐◌◌│◌║ 4} \\
\texttt{║◌◌├──┼─╢ 5} \\
\texttt{║◌W│◌◌│C║ 6} \\
\texttt{║◌◌│◌◌│◌║ 7} \\
\texttt{╚══╧══╧═╝} \\
\texttt{```} \\
\texttt{Tá sibh ag imirt babhta 5 faoi láthair as uasmhéid 20.} \\
\texttt{Tá 2 as 12 pionós carntha agaibh le chéile.} \\
\texttt{Sheol an t{-}imreoir eile an teachtaireacht seo:} \\
\\ 
\\ 
\texttt{"D'éirigh liom T a bhogadh go (1,1). Ach tá frámaí ag (6,6) agus (7,7) ar mo chlár, mar sin ní féidir liom spriocanna a gcuid a shroichint ansin. Molaim sprioc nua: I ag (2,6), C ag (3,7), H ag (4,6). Cad é do bharúil?"} \\
\\ 
\\ 
\texttt{Cad é do chéad ordú eile?} \\
            }
        }
    }
    & \\ \\

    \theutterance \stepcounter{utterance}  
    & & & \multicolumn{4}{p{0.6\linewidth}}{
        \cellcolor[rgb]{0.9,0.9,0.9}{
            \makecell[{{p{\linewidth}}}]{
                \texttt{\tiny{[GM$\langle$P2]}}
                \texttt{ABAIR: Tá brón orm, ní féidir liom spriocanna a ghlacadh ag (2,6) mar go bhfuil W cheana ann, ná ag (3,7) mar gheall ar fhráma agus (4,6) as theorainn. Molaim sprioc úr atá saor ó fhrámaí ar mo chlár: W ag (2,4), T atá agat ag (1,1), I ag (3,1), C atá ag (1,7), H ag (4,5). Conas é sin?} \\
            }
        }
    }
     \\ \\

    \theutterance \stepcounter{utterance}  
    & & & \multicolumn{2}{p{0.3\linewidth}}{
        \cellcolor[rgb]{0.9,0.9,0.9}{
            \makecell[{{p{\linewidth}}}]{
                \texttt{\tiny{[GM$|$GM]}}
                \texttt{ABAIR: Tá brón orm, ní féidir liom spriocanna a ghlacadh ag (2,6) mar go bhfuil W cheana ann, ná ag (3,7) mar gheall ar fhráma agus (4,6) as theorainn. Molaim sprioc úr atá saor ó fhrámaí ar mo chlár: W ag (2,4), T atá agat ag (1,1), I ag (3,1), C atá ag (1,7), H ag (4,5). Conas é sin?} \\
            }
        }
    }
    & & \\ \\

    \theutterance \stepcounter{utterance}  
    & & \multicolumn{4}{p{0.6\linewidth}}{
        \cellcolor[rgb]{0.9,0.9,0.9}{
            \makecell[{{p{\linewidth}}}]{
                \texttt{\tiny{[P1$\langle$GM]}}
                \texttt{Cuireadh do theachtaireacht ar aghaidh chuig an imreoir eile.} \\
\\ 
\\ 
\texttt{Tá sibh ag imirt babhta 5 faoi láthair as uasmhéid 20.} \\
\texttt{Tá 2 as 12 pionós carntha agaibh le chéile.} \\
\texttt{Sheol an t{-}imreoir eile an teachtaireacht seo:} \\
\\ 
\\ 
\texttt{"Tá brón orm, ní féidir liom spriocanna a ghlacadh ag (2,6) mar go bhfuil W cheana ann, ná ag (3,7) mar gheall ar fhráma agus (4,6) as theorainn. Molaim sprioc úr atá saor ó fhrámaí ar mo chlár: W ag (2,4), T atá agat ag (1,1), I ag (3,1), C atá ag (1,7), H ag (4,5). Conas é sin?"} \\
\\ 
\\ 
\texttt{Cad é do chéad ordú eile?} \\
            }
        }
    }
    & \\ \\

    \theutterance \stepcounter{utterance}  
    & \multicolumn{4}{p{0.6\linewidth}}{
        \cellcolor[rgb]{0.9,0.9,0.9}{
            \makecell[{{p{\linewidth}}}]{
                \texttt{\tiny{[P1$\rangle$GM]}}
                \texttt{ABAIR: Tá fráma ag (4,5) ar mo chlár freisin. Molaim spriocanna i láithreacha folmhacha cinnte ar an dá chlár: W ag (5,1), I ag (5,2), C ag (5,3), H ag (6,7). Fan go dtiocfaidh an dá thaobh ar aon fhocal sula ndéanaimid bogadh.} \\
            }
        }
    }
    & & \\ \\

    \theutterance \stepcounter{utterance}  
    & & & \multicolumn{2}{p{0.3\linewidth}}{
        \cellcolor[rgb]{0.9,0.9,0.9}{
            \makecell[{{p{\linewidth}}}]{
                \texttt{\tiny{[GM$|$GM]}}
                \texttt{ABAIR: Tá fráma ag (4,5) ar mo chlár freisin. Molaim spriocanna i láithreacha folmhacha cinnte ar an dá chlár: W ag (5,1), I ag (5,2), C ag (5,3), H ag (6,7). Fan go dtiocfaidh an dá thaobh ar aon fhocal sula ndéanaimid bogadh.} \\
            }
        }
    }
    & & \\ \\

    \theutterance \stepcounter{utterance}  
    & & \multicolumn{4}{p{0.6\linewidth}}{
        \cellcolor[rgb]{0.9,0.9,0.9}{
            \makecell[{{p{\linewidth}}}]{
                \texttt{\tiny{[GM$\rangle$P2]}}
                \texttt{Cuireadh do theachtaireacht ar aghaidh chuig an imreoir eile.} \\
\\ 
\\ 
\texttt{Tá sibh ag imirt babhta 6 faoi láthair as uasmhéid 20.} \\
\texttt{Tá 2 as 12 pionós carntha agaibh le chéile.} \\
\texttt{Sheol an t{-}imreoir eile an teachtaireacht seo:} \\
\\ 
\\ 
\texttt{"Tá fráma ag (4,5) ar mo chlár freisin. Molaim spriocanna i láithreacha folmhacha cinnte ar an dá chlár: W ag (5,1), I ag (5,2), C ag (5,3), H ag (6,7). Fan go dtiocfaidh an dá thaobh ar aon fhocal sula ndéanaimid bogadh."} \\
\\ 
\\ 
\texttt{Cad é do chéad ordú eile?} \\
            }
        }
    }
    & \\ \\

    \theutterance \stepcounter{utterance}  
    & & & \multicolumn{4}{p{0.6\linewidth}}{
        \cellcolor[rgb]{0.9,0.9,0.9}{
            \makecell[{{p{\linewidth}}}]{
                \texttt{\tiny{[GM$\langle$P2]}}
                \texttt{ABAIR: Tá mo chlár an{-}chosúil le do cheann ar feadh na líne seo! Tá fráma agam freisin ag (5,3), áit ar mhaith leat C a chur. Mar sin féin, tá na cealla eile go léir folamh agam. Tá an plean seo go hiomlán míléiteartha ó thaobh na loighnice greille de: W ag (4,7), I ag (5,7), C atá agat (1,1), H ag (2,2), T ag (6,6). Glacaimid leis an socrú seo?} \\
            }
        }
    }
     \\ \\

    \theutterance \stepcounter{utterance}  
    & & & \multicolumn{2}{p{0.3\linewidth}}{
        \cellcolor[rgb]{0.9,0.9,0.9}{
            \makecell[{{p{\linewidth}}}]{
                \texttt{\tiny{[GM$|$GM]}}
                \texttt{ABAIR: Tá mo chlár an{-}chosúil le do cheann ar feadh na líne seo! Tá fráma agam freisin ag (5,3), áit ar mhaith leat C a chur. Mar sin féin, tá na cealla eile go léir folamh agam. Tá an plean seo go hiomlán míléiteartha ó thaobh na loighnice greille de: W ag (4,7), I ag (5,7), C atá agat (1,1), H ag (2,2), T ag (6,6). Glacaimid leis an socrú seo?} \\
            }
        }
    }
    & & \\ \\

    \theutterance \stepcounter{utterance}  
    & & \multicolumn{4}{p{0.6\linewidth}}{
        \cellcolor[rgb]{0.9,0.9,0.9}{
            \makecell[{{p{\linewidth}}}]{
                \texttt{\tiny{[P1$\langle$GM]}}
                \texttt{Cuireadh do theachtaireacht ar aghaidh chuig an imreoir eile.} \\
\\ 
\\ 
\texttt{Tá sibh ag imirt babhta 6 faoi láthair as uasmhéid 20.} \\
\texttt{Tá 2 as 12 pionós carntha agaibh le chéile.} \\
\texttt{Sheol an t{-}imreoir eile an teachtaireacht seo:} \\
\\ 
\\ 
\texttt{"Tá mo chlár an{-}chosúil le do cheann ar feadh na líne seo! Tá fráma agam freisin ag (5,3), áit ar mhaith leat C a chur. Mar sin féin, tá na cealla eile go léir folamh agam. Tá an plean seo go hiomlán míléiteartha ó thaobh na loighnice greille de: W ag (4,7), I ag (5,7), C atá agat (1,1), H ag (2,2), T ag (6,6). Glacaimid leis an socrú seo?"} \\
\\ 
\\ 
\texttt{Cad é do chéad ordú eile?} \\
            }
        }
    }
    & \\ \\

    \theutterance \stepcounter{utterance}  
    & \multicolumn{4}{p{0.6\linewidth}}{
        \cellcolor[rgb]{0.9,0.9,0.9}{
            \makecell[{{p{\linewidth}}}]{
                \texttt{\tiny{[P1$\rangle$GM]}}
                \texttt{ABAIR: Tá fráma ag (6,6) ar mo chlár freisin! Caithfimid sprioc a fháil gan frámaí ar chor ar bith. Seo iad na spotaí folamh glan atá agam: (1,2), (2,2), (3,2), (4,2), (6,2), (7,2), (1,6), (2,7), (3,7), (4,6), (6,7), (7,4), (7,5). An bhfuil aon cheann de seo folamh glan agatsa freisin? Roghnóidh mé as iad sin.} \\
            }
        }
    }
    & & \\ \\

    \theutterance \stepcounter{utterance}  
    & & & \multicolumn{2}{p{0.3\linewidth}}{
        \cellcolor[rgb]{0.9,0.9,0.9}{
            \makecell[{{p{\linewidth}}}]{
                \texttt{\tiny{[GM$|$GM]}}
                \texttt{ABAIR: Tá fráma ag (6,6) ar mo chlár freisin! Caithfimid sprioc a fháil gan frámaí ar chor ar bith. Seo iad na spotaí folamh glan atá agam: (1,2), (2,2), (3,2), (4,2), (6,2), (7,2), (1,6), (2,7), (3,7), (4,6), (6,7), (7,4), (7,5). An bhfuil aon cheann de seo folamh glan agatsa freisin? Roghnóidh mé as iad sin.} \\
            }
        }
    }
    & & \\ \\

    \theutterance \stepcounter{utterance}  
    & & \multicolumn{4}{p{0.6\linewidth}}{
        \cellcolor[rgb]{0.9,0.9,0.9}{
            \makecell[{{p{\linewidth}}}]{
                \texttt{\tiny{[GM$\rangle$P2]}}
                \texttt{Cuireadh do theachtaireacht ar aghaidh chuig an imreoir eile.} \\
\\ 
\\ 
\texttt{Tá sibh ag imirt babhta 7 faoi láthair as uasmhéid 20.} \\
\texttt{Tá 2 as 12 pionós carntha agaibh le chéile.} \\
\texttt{Sheol an t{-}imreoir eile an teachtaireacht seo:} \\
\\ 
\\ 
\texttt{"Tá fráma ag (6,6) ar mo chlár freisin! Caithfimid sprioc a fháil gan frámaí ar chor ar bith. Seo iad na spotaí folamh glan atá agam: (1,2), (2,2), (3,2), (4,2), (6,2), (7,2), (1,6), (2,7), (3,7), (4,6), (6,7), (7,4), (7,5). An bhfuil aon cheann de seo folamh glan agatsa freisin? Roghnóidh mé as iad sin."} \\
\\ 
\\ 
\texttt{Cad é do chéad ordú eile?} \\
            }
        }
    }
    & \\ \\

    \theutterance \stepcounter{utterance}  
    & & & \multicolumn{4}{p{0.6\linewidth}}{
        \cellcolor[rgb]{0.9,0.9,0.9}{
            \makecell[{{p{\linewidth}}}]{
                \texttt{\tiny{[GM$\langle$P2]}}
                \texttt{ABAIR: Breathnaigh, níl gá le hairdre focal! Is féidir linn dul i dtaithí ar an tsiombailceann agus an bheirt a bheith aontaithe. Seo na hoibríochtaí go léir is féidir a dhéanamh leis na ballaí atá ar fud an bhaile. Seiceáil na horduithe seo a leanas, agus oscailfidh tú doras nua:} \\
\texttt{Tá na frámaí ar do chlár níos casta ná mar a shíl mé. Mar sin, caithfimid bheith cúramach.} \\
\texttt{Seo iad na spotaí folamh glan atá agamsa: (1,1), (2,1), (3,1), (4,1), (5,1), (6,1), (7,1), (3,2), (4,2), (5,2), (7,2), (4,4), (7,4), (5,5), (7,5), (1,6), (3,6), (4,6), (6,7). Tá (3,7) in aice leis an bhfráma ach glan.} \\
\texttt{Is iad seo na hionaid a mholaim: W ag (7,1), T ag (1,1), H ag (7,2), I ag (5,2), C ag (7,5). Roghnóidh mé as na spotaí atá folamh glan agamsa freisin. An bhfuil an plean seo go maith?} \\
            }
        }
    }
     \\ \\

    \theutterance \stepcounter{utterance}  
    & & & \multicolumn{2}{p{0.3\linewidth}}{
        \cellcolor[rgb]{0.9,0.9,0.9}{
            \makecell[{{p{\linewidth}}}]{
                \texttt{\tiny{[GM$|$GM]}}
                \texttt{ABAIR: Breathnaigh, níl gá le hairdre focal! Is féidir linn dul i dtaithí ar an tsiombailceann agus an bheirt a bheith aontaithe. Seo na hoibríochtaí go léir is féidir a dhéanamh leis na ballaí atá ar fud an bhaile. Seiceáil na horduithe seo a leanas, agus oscailfidh tú doras nua:} \\
\texttt{Tá na frámaí ar do chlár níos casta ná mar a shíl mé. Mar sin, caithfimid bheith cúramach.} \\
\texttt{Seo iad na spotaí folamh glan atá agamsa: (1,1), (2,1), (3,1), (4,1), (5,1), (6,1), (7,1), (3,2), (4,2), (5,2), (7,2), (4,4), (7,4), (5,5), (7,5), (1,6), (3,6), (4,6), (6,7). Tá (3,7) in aice leis an bhfráma ach glan.} \\
\texttt{Is iad seo na hionaid a mholaim: W ag (7,1), T ag (1,1), H ag (7,2), I ag (5,2), C ag (7,5). Roghnóidh mé as na spotaí atá folamh glan agamsa freisin. An bhfuil an plean seo go maith?} \\
            }
        }
    }
    & & \\ \\

    \theutterance \stepcounter{utterance}  
    & & & \multicolumn{2}{p{0.3\linewidth}}{
        \cellcolor[rgb]{0.9,0.9,0.9}{
            \makecell[{{p{\linewidth}}}]{
                \texttt{\tiny{[GM$|$GM]}}
                \texttt{Invalid format: tail is not empty: '} \\
\texttt{Tá na frámaí ar do chlár níos casta ná mar a shíl mé. Mar sin, caithfimid bheith cúramach.} \\
\texttt{Seo iad na spotaí folamh glan atá agamsa: (1,1), (2,1), (3,1), (4,1), (5,1), (6,1), (7,1), (3,2), (4,2), (5,2), (7,2), (4,4), (7,4), (5,5), (7,5), (1,6), (3,6), (4,6), (6,7). Tá (3,7) in aice leis an bhfráma ach glan.} \\
\texttt{Is iad seo na hionaid a mholaim: W ag (7,1), T ag (1,1), H ag (7,2), I ag (5,2), C ag (7,5). Roghnóidh mé as na spotaí atá folamh glan agamsa freisin. An bhfuil an plean seo go maith?} \\
            }
        }
    }
    & & \\ \\

    \theutterance \stepcounter{utterance}  
    & & \multicolumn{4}{p{0.6\linewidth}}{
        \cellcolor[rgb]{0.9,0.9,0.9}{
            \makecell[{{p{\linewidth}}}]{
                \texttt{\tiny{[GM$\rangle$P2]}}
                \texttt{Pionós: Ní cheadaítear aon rud a bheith i do theachtaireacht tar éis an ordaithe!} \\
\\ 
\texttt{Cinntigh nach bhfuil i do fhreagra ach `ABAIR: <TEACHTAIREACHT>` nó `BOG: <RUD>, (<X>, <Y>)`, agus tada eile!} \\
\\ 
\texttt{Tá 3 as 12 pionós carntha agaibh le chéile.} \\
\texttt{Déan iarracht arís, le do thoil!} \\
            }
        }
    }
    & \\ \\

    \theutterance \stepcounter{utterance}  
    & & & \multicolumn{4}{p{0.6\linewidth}}{
        \cellcolor[rgb]{0.9,0.9,0.9}{
            \makecell[{{p{\linewidth}}}]{
                \texttt{\tiny{[GM$\langle$P2]}}
                \texttt{ABAIR: Tá na spotaí folamh glan atá agamsa: (1,1), (2,1), (3,1), (4,1), (5,1), (6,1), (7,1), (3,2), (4,2), (5,2), (7,2), (4,4), (7,4), (5,5), (7,5), (1,6), (3,6), (4,6), (6,7). Is iad seo na hionaid a mholaim: W ag (7,1), T ag (1,1), H ag (7,2), I ag (5,2), C ag (7,5). Roghnóidh mé as na spotaí atá folamh glan agamsa freisin. An bhfuil an plean seo go maith?} \\
            }
        }
    }
     \\ \\

    \theutterance \stepcounter{utterance}  
    & & & \multicolumn{2}{p{0.3\linewidth}}{
        \cellcolor[rgb]{0.9,0.9,0.9}{
            \makecell[{{p{\linewidth}}}]{
                \texttt{\tiny{[GM$|$GM]}}
                \texttt{ABAIR: Tá na spotaí folamh glan atá agamsa: (1,1), (2,1), (3,1), (4,1), (5,1), (6,1), (7,1), (3,2), (4,2), (5,2), (7,2), (4,4), (7,4), (5,5), (7,5), (1,6), (3,6), (4,6), (6,7). Is iad seo na hionaid a mholaim: W ag (7,1), T ag (1,1), H ag (7,2), I ag (5,2), C ag (7,5). Roghnóidh mé as na spotaí atá folamh glan agamsa freisin. An bhfuil an plean seo go maith?} \\
            }
        }
    }
    & & \\ \\

    \theutterance \stepcounter{utterance}  
    & & \multicolumn{4}{p{0.6\linewidth}}{
        \cellcolor[rgb]{0.9,0.9,0.9}{
            \makecell[{{p{\linewidth}}}]{
                \texttt{\tiny{[P1$\langle$GM]}}
                \texttt{Cuireadh do theachtaireacht ar aghaidh chuig an imreoir eile.} \\
\\ 
\\ 
\texttt{Tá sibh ag imirt babhta 7 faoi láthair as uasmhéid 20.} \\
\texttt{Tá 3 as 12 pionós carntha agaibh le chéile.} \\
\texttt{Sheol an t{-}imreoir eile an teachtaireacht seo:} \\
\\ 
\\ 
\texttt{"Tá na spotaí folamh glan atá agamsa: (1,1), (2,1), (3,1), (4,1), (5,1), (6,1), (7,1), (3,2), (4,2), (5,2), (7,2), (4,4), (7,4), (5,5), (7,5), (1,6), (3,6), (4,6), (6,7). Is iad seo na hionaid a mholaim: W ag (7,1), T ag (1,1), H ag (7,2), I ag (5,2), C ag (7,5). Roghnóidh mé as na spotaí atá folamh glan agamsa freisin. An bhfuil an plean seo go maith?"} \\
\\ 
\\ 
\texttt{Cad é do chéad ordú eile?} \\
            }
        }
    }
    & \\ \\

    \theutterance \stepcounter{utterance}  
    & \multicolumn{4}{p{0.6\linewidth}}{
        \cellcolor[rgb]{0.9,0.9,0.9}{
            \makecell[{{p{\linewidth}}}]{
                \texttt{\tiny{[P1$\rangle$GM]}}
                \texttt{ABAIR: Tá T agam cheana féin ag (1,1), is maith sin. Seiceáil mo shliotáin folamh glan eile atá ar fhíor{-}spás: (1,2), (2,2), (3,2), (4,2), (6,2), (7,2), (1,6), (4,6), (6,7). Mholaim: W ag (7,2), H ag (1,2), I ag (6,2), C ag (6,7). Ceadaigh?} \\
            }
        }
    }
    & & \\ \\

    \theutterance \stepcounter{utterance}  
    & & & \multicolumn{2}{p{0.3\linewidth}}{
        \cellcolor[rgb]{0.9,0.9,0.9}{
            \makecell[{{p{\linewidth}}}]{
                \texttt{\tiny{[GM$|$GM]}}
                \texttt{ABAIR: Tá T agam cheana féin ag (1,1), is maith sin. Seiceáil mo shliotáin folamh glan eile atá ar fhíor{-}spás: (1,2), (2,2), (3,2), (4,2), (6,2), (7,2), (1,6), (4,6), (6,7). Mholaim: W ag (7,2), H ag (1,2), I ag (6,2), C ag (6,7). Ceadaigh?} \\
            }
        }
    }
    & & \\ \\

    \theutterance \stepcounter{utterance}  
    & & \multicolumn{4}{p{0.6\linewidth}}{
        \cellcolor[rgb]{0.9,0.9,0.9}{
            \makecell[{{p{\linewidth}}}]{
                \texttt{\tiny{[GM$\rangle$P2]}}
                \texttt{Cuireadh do theachtaireacht ar aghaidh chuig an imreoir eile.} \\
\\ 
\\ 
\texttt{Tá sibh ag imirt babhta 8 faoi láthair as uasmhéid 20.} \\
\texttt{Tá 3 as 12 pionós carntha agaibh le chéile.} \\
\texttt{Sheol an t{-}imreoir eile an teachtaireacht seo:} \\
\\ 
\\ 
\texttt{"Tá T agam cheana féin ag (1,1), is maith sin. Seiceáil mo shliotáin folamh glan eile atá ar fhíor{-}spás: (1,2), (2,2), (3,2), (4,2), (6,2), (7,2), (1,6), (4,6), (6,7). Mholaim: W ag (7,2), H ag (1,2), I ag (6,2), C ag (6,7). Ceadaigh?"} \\
\\ 
\\ 
\texttt{Cad é do chéad ordú eile?} \\
            }
        }
    }
    & \\ \\

    \theutterance \stepcounter{utterance}  
    & & & \multicolumn{4}{p{0.6\linewidth}}{
        \cellcolor[rgb]{0.9,0.9,0.9}{
            \makecell[{{p{\linewidth}}}]{
                \texttt{\tiny{[GM$\langle$P2]}}
                \texttt{ABAIR: Tá (6,7) glan agam ach tá (6,2) agus (7,2) agamsa le frámaí nó bac. Molaim sprioc ina ionad: W ag (4,1), H ag (2,1), I ag (3,2), C ag (6,7). An bhfuil sin ceart agat?} \\
            }
        }
    }
     \\ \\

    \theutterance \stepcounter{utterance}  
    & & & \multicolumn{2}{p{0.3\linewidth}}{
        \cellcolor[rgb]{0.9,0.9,0.9}{
            \makecell[{{p{\linewidth}}}]{
                \texttt{\tiny{[GM$|$GM]}}
                \texttt{ABAIR: Tá (6,7) glan agam ach tá (6,2) agus (7,2) agamsa le frámaí nó bac. Molaim sprioc ina ionad: W ag (4,1), H ag (2,1), I ag (3,2), C ag (6,7). An bhfuil sin ceart agat?} \\
            }
        }
    }
    & & \\ \\

    \theutterance \stepcounter{utterance}  
    & & \multicolumn{4}{p{0.6\linewidth}}{
        \cellcolor[rgb]{0.9,0.9,0.9}{
            \makecell[{{p{\linewidth}}}]{
                \texttt{\tiny{[P1$\langle$GM]}}
                \texttt{Cuireadh do theachtaireacht ar aghaidh chuig an imreoir eile.} \\
\\ 
\\ 
\texttt{Tá sibh ag imirt babhta 8 faoi láthair as uasmhéid 20.} \\
\texttt{Tá 3 as 12 pionós carntha agaibh le chéile.} \\
\texttt{Sheol an t{-}imreoir eile an teachtaireacht seo:} \\
\\ 
\\ 
\texttt{"Tá (6,7) glan agam ach tá (6,2) agus (7,2) agamsa le frámaí nó bac. Molaim sprioc ina ionad: W ag (4,1), H ag (2,1), I ag (3,2), C ag (6,7). An bhfuil sin ceart agat?"} \\
\\ 
\\ 
\texttt{Cad é do chéad ordú eile?} \\
            }
        }
    }
    & \\ \\

    \theutterance \stepcounter{utterance}  
    & \multicolumn{4}{p{0.6\linewidth}}{
        \cellcolor[rgb]{0.9,0.9,0.9}{
            \makecell[{{p{\linewidth}}}]{
                \texttt{\tiny{[P1$\rangle$GM]}}
                \texttt{ABAIR: Tá (4,1), (2,1), (3,2) agus (6,7) go léir folamh glan agam freisin! Aontaím go huile is go hiomlán leis an sprioc nua seo. Tosóidh mé ag bogadh W go (4,1) anois. Bog saoirse H go (2,1) agus inis dom céard atá fágaithe agat.} \\
            }
        }
    }
    & & \\ \\

    \theutterance \stepcounter{utterance}  
    & & & \multicolumn{2}{p{0.3\linewidth}}{
        \cellcolor[rgb]{0.9,0.9,0.9}{
            \makecell[{{p{\linewidth}}}]{
                \texttt{\tiny{[GM$|$GM]}}
                \texttt{ABAIR: Tá (4,1), (2,1), (3,2) agus (6,7) go léir folamh glan agam freisin! Aontaím go huile is go hiomlán leis an sprioc nua seo. Tosóidh mé ag bogadh W go (4,1) anois. Bog saoirse H go (2,1) agus inis dom céard atá fágaithe agat.} \\
            }
        }
    }
    & & \\ \\

    \theutterance \stepcounter{utterance}  
    & & \multicolumn{4}{p{0.6\linewidth}}{
        \cellcolor[rgb]{0.9,0.9,0.9}{
            \makecell[{{p{\linewidth}}}]{
                \texttt{\tiny{[GM$\rangle$P2]}}
                \texttt{Cuireadh do theachtaireacht ar aghaidh chuig an imreoir eile.} \\
\\ 
\\ 
\texttt{Tá sibh ag imirt babhta 9 faoi láthair as uasmhéid 20.} \\
\texttt{Tá 3 as 12 pionós carntha agaibh le chéile.} \\
\texttt{Sheol an t{-}imreoir eile an teachtaireacht seo:} \\
\\ 
\\ 
\texttt{"Tá (4,1), (2,1), (3,2) agus (6,7) go léir folamh glan agam freisin! Aontaím go huile is go hiomlán leis an sprioc nua seo. Tosóidh mé ag bogadh W go (4,1) anois. Bog saoirse H go (2,1) agus inis dom céard atá fágaithe agat."} \\
\\ 
\\ 
\texttt{Cad é do chéad ordú eile?} \\
            }
        }
    }
    & \\ \\

    \theutterance \stepcounter{utterance}  
    & & & \multicolumn{4}{p{0.6\linewidth}}{
        \cellcolor[rgb]{0.9,0.9,0.9}{
            \makecell[{{p{\linewidth}}}]{
                \texttt{\tiny{[GM$\langle$P2]}}
                \texttt{BOG: I, (3,2)} \\
            }
        }
    }
     \\ \\

    \theutterance \stepcounter{utterance}  
    & & & \multicolumn{2}{p{0.3\linewidth}}{
        \cellcolor[rgb]{0.9,0.9,0.9}{
            \makecell[{{p{\linewidth}}}]{
                \texttt{\tiny{[GM$|$GM]}}
                \texttt{BOG: I, (3,2)} \\
            }
        }
    }
    & & \\ \\

    \theutterance \stepcounter{utterance}  
    & & & \multicolumn{2}{p{0.3\linewidth}}{
        \cellcolor[rgb]{0.9,0.9,0.9}{
            \makecell[{{p{\linewidth}}}]{
                \texttt{\tiny{[GM$|$GM]}}
                \texttt{I moved to (3, 2)} \\
            }
        }
    }
    & & \\ \\

    \theutterance \stepcounter{utterance}  
    & & \multicolumn{4}{p{0.6\linewidth}}{
        \cellcolor[rgb]{0.9,0.9,0.9}{
            \makecell[{{p{\linewidth}}}]{
                \texttt{\tiny{[P1$\langle$GM]}}
                \texttt{Cuireadh do theachtaireacht ar aghaidh chuig an imreoir eile.} \\
\\ 
\\ 
\texttt{Tá sibh ag imirt babhta 9 faoi láthair as uasmhéid 20.} \\
\texttt{Tá 3 as 12 pionós carntha agaibh le chéile.} \\
\texttt{Bhog an t{-}imreoir eile réad ar a ngreille.} \\
\\ 
\\ 
\texttt{Cad é do chéad ordú eile?} \\
            }
        }
    }
    & \\ \\

    \theutterance \stepcounter{utterance}  
    & \multicolumn{4}{p{0.6\linewidth}}{
        \cellcolor[rgb]{0.9,0.9,0.9}{
            \makecell[{{p{\linewidth}}}]{
                \texttt{\tiny{[P1$\rangle$GM]}}
                \texttt{BOG: W, (4, 1)} \\
            }
        }
    }
    & & \\ \\

    \theutterance \stepcounter{utterance}  
    & & & \multicolumn{2}{p{0.3\linewidth}}{
        \cellcolor[rgb]{0.9,0.9,0.9}{
            \makecell[{{p{\linewidth}}}]{
                \texttt{\tiny{[GM$|$GM]}}
                \texttt{BOG: W, (4, 1)} \\
            }
        }
    }
    & & \\ \\

    \theutterance \stepcounter{utterance}  
    & & & \multicolumn{2}{p{0.3\linewidth}}{
        \cellcolor[rgb]{0.9,0.9,0.9}{
            \makecell[{{p{\linewidth}}}]{
                \texttt{\tiny{[GM$|$GM]}}
                \texttt{W moved to (4, 1)} \\
            }
        }
    }
    & & \\ \\

    \theutterance \stepcounter{utterance}  
    & & \multicolumn{4}{p{0.6\linewidth}}{
        \cellcolor[rgb]{0.9,0.9,0.9}{
            \makecell[{{p{\linewidth}}}]{
                \texttt{\tiny{[GM$\rangle$P2]}}
                \texttt{Bogadh 'I' go (3,2) go rathúil. Seo mar atá do ghreille nuashonraithe:} \\
\texttt{```} \\
\texttt{ 1234567} \\
\texttt{╔═══════╗} \\
\texttt{║◌◌◌◌TH◌║ 1} \\
\texttt{║◌◌I◌◌◌◌║ 2} \\
\texttt{║◌◌◌◌◌┌─╢ 3} \\
\texttt{╟──┐◌◌│◌║ 4} \\
\texttt{║◌◌├──┼─╢ 5} \\
\texttt{║◌W│◌◌│C║ 6} \\
\texttt{║◌◌│◌◌│◌║ 7} \\
\texttt{╚══╧══╧═╝} \\
\texttt{```} \\
\texttt{Tá sibh ag imirt babhta 10 faoi láthair as uasmhéid 20.} \\
\texttt{Tá 3 as 12 pionós carntha agaibh le chéile.} \\
\texttt{Bhog an t{-}imreoir eile réad ar a ngreille.} \\
\\ 
\\ 
\texttt{Cad é do chéad ordú eile?} \\
            }
        }
    }
    & \\ \\

    \theutterance \stepcounter{utterance}  
    & & & \multicolumn{4}{p{0.6\linewidth}}{
        \cellcolor[rgb]{0.9,0.9,0.9}{
            \makecell[{{p{\linewidth}}}]{
                \texttt{\tiny{[GM$\langle$P2]}}
                \texttt{ABAIR: D'éirigh liom I a bhogadh go (3,2). Anois bogfaidh mé H go (2,1). Inis dom conas a d'éirigh le do bhogadh W.} \\
            }
        }
    }
     \\ \\

    \theutterance \stepcounter{utterance}  
    & & & \multicolumn{2}{p{0.3\linewidth}}{
        \cellcolor[rgb]{0.9,0.9,0.9}{
            \makecell[{{p{\linewidth}}}]{
                \texttt{\tiny{[GM$|$GM]}}
                \texttt{ABAIR: D'éirigh liom I a bhogadh go (3,2). Anois bogfaidh mé H go (2,1). Inis dom conas a d'éirigh le do bhogadh W.} \\
            }
        }
    }
    & & \\ \\

    \theutterance \stepcounter{utterance}  
    & & \multicolumn{4}{p{0.6\linewidth}}{
        \cellcolor[rgb]{0.9,0.9,0.9}{
            \makecell[{{p{\linewidth}}}]{
                \texttt{\tiny{[P1$\langle$GM]}}
                \texttt{Bogadh 'W' go (4,1) go rathúil. Seo mar atá do ghreille nuashonraithe:} \\
\texttt{```} \\
\texttt{ 1234567} \\
\texttt{╔═══════╗} \\
\texttt{║T◌◌W◌◌◌║ 1} \\
\texttt{║◌◌◌◌◌◌◌║ 2} \\
\texttt{║◌◌◌◌◌┌─╢ 3} \\
\texttt{╟──┐◌◌│H║ 4} \\
\texttt{║C◌├──┼─╢ 5} \\
\texttt{║◌◌│I◌│◌║ 6} \\
\texttt{║◌◌│◌◌│◌║ 7} \\
\texttt{╚══╧══╧═╝} \\
\texttt{```} \\
\texttt{Tá sibh ag imirt babhta 10 faoi láthair as uasmhéid 20.} \\
\texttt{Tá 3 as 12 pionós carntha agaibh le chéile.} \\
\texttt{Sheol an t{-}imreoir eile an teachtaireacht seo:} \\
\\ 
\\ 
\texttt{"D'éirigh liom I a bhogadh go (3,2). Anois bogfaidh mé H go (2,1). Inis dom conas a d'éirigh le do bhogadh W."} \\
\\ 
\\ 
\texttt{Cad é do chéad ordú eile?} \\
            }
        }
    }
    & \\ \\

    \theutterance \stepcounter{utterance}  
    & \multicolumn{4}{p{0.6\linewidth}}{
        \cellcolor[rgb]{0.9,0.9,0.9}{
            \makecell[{{p{\linewidth}}}]{
                \texttt{\tiny{[P1$\rangle$GM]}}
                \texttt{ABAIR: D'éirigh go hálainn liom W a bhogadh go (4,1). Tá T ag (1,1), W ag (4,1), agus tú ag bogadh I go (3,2) agus H go (2,1). Ní fágfaidh sin ach C agam le bogadh go (6,7). Bogfaidh mé C i mo chéad chéad seans eile. Conas a d'éirigh leat le H?} \\
            }
        }
    }
    & & \\ \\

    \theutterance \stepcounter{utterance}  
    & & & \multicolumn{2}{p{0.3\linewidth}}{
        \cellcolor[rgb]{0.9,0.9,0.9}{
            \makecell[{{p{\linewidth}}}]{
                \texttt{\tiny{[GM$|$GM]}}
                \texttt{ABAIR: D'éirigh go hálainn liom W a bhogadh go (4,1). Tá T ag (1,1), W ag (4,1), agus tú ag bogadh I go (3,2) agus H go (2,1). Ní fágfaidh sin ach C agam le bogadh go (6,7). Bogfaidh mé C i mo chéad chéad seans eile. Conas a d'éirigh leat le H?} \\
            }
        }
    }
    & & \\ \\

    \theutterance \stepcounter{utterance}  
    & & \multicolumn{4}{p{0.6\linewidth}}{
        \cellcolor[rgb]{0.9,0.9,0.9}{
            \makecell[{{p{\linewidth}}}]{
                \texttt{\tiny{[GM$\rangle$P2]}}
                \texttt{Cuireadh do theachtaireacht ar aghaidh chuig an imreoir eile.} \\
\\ 
\\ 
\texttt{Tá sibh ag imirt babhta 11 faoi láthair as uasmhéid 20.} \\
\texttt{Tá 3 as 12 pionós carntha agaibh le chéile.} \\
\texttt{Sheol an t{-}imreoir eile an teachtaireacht seo:} \\
\\ 
\\ 
\texttt{"D'éirigh go hálainn liom W a bhogadh go (4,1). Tá T ag (1,1), W ag (4,1), agus tú ag bogadh I go (3,2) agus H go (2,1). Ní fágfaidh sin ach C agam le bogadh go (6,7). Bogfaidh mé C i mo chéad chéad seans eile. Conas a d'éirigh leat le H?"} \\
\\ 
\\ 
\texttt{Cad é do chéad ordú eile?} \\
            }
        }
    }
    & \\ \\

    \theutterance \stepcounter{utterance}  
    & & & \multicolumn{4}{p{0.6\linewidth}}{
        \cellcolor[rgb]{0.9,0.9,0.9}{
            \makecell[{{p{\linewidth}}}]{
                \texttt{\tiny{[GM$\langle$P2]}}
                \texttt{BOG: H, (2, 1)} \\
            }
        }
    }
     \\ \\

    \theutterance \stepcounter{utterance}  
    & & & \multicolumn{2}{p{0.3\linewidth}}{
        \cellcolor[rgb]{0.9,0.9,0.9}{
            \makecell[{{p{\linewidth}}}]{
                \texttt{\tiny{[GM$|$GM]}}
                \texttt{BOG: H, (2, 1)} \\
            }
        }
    }
    & & \\ \\

    \theutterance \stepcounter{utterance}  
    & & & \multicolumn{2}{p{0.3\linewidth}}{
        \cellcolor[rgb]{0.9,0.9,0.9}{
            \makecell[{{p{\linewidth}}}]{
                \texttt{\tiny{[GM$|$GM]}}
                \texttt{H moved to (2, 1)} \\
            }
        }
    }
    & & \\ \\

    \theutterance \stepcounter{utterance}  
    & & \multicolumn{4}{p{0.6\linewidth}}{
        \cellcolor[rgb]{0.9,0.9,0.9}{
            \makecell[{{p{\linewidth}}}]{
                \texttt{\tiny{[P1$\langle$GM]}}
                \texttt{Cuireadh do theachtaireacht ar aghaidh chuig an imreoir eile.} \\
\\ 
\\ 
\texttt{Tá sibh ag imirt babhta 11 faoi láthair as uasmhéid 20.} \\
\texttt{Tá 3 as 12 pionós carntha agaibh le chéile.} \\
\texttt{Bhog an t{-}imreoir eile réad ar a ngreille.} \\
\\ 
\\ 
\texttt{Cad é do chéad ordú eile?} \\
            }
        }
    }
    & \\ \\

    \theutterance \stepcounter{utterance}  
    & \multicolumn{4}{p{0.6\linewidth}}{
        \cellcolor[rgb]{0.9,0.9,0.9}{
            \makecell[{{p{\linewidth}}}]{
                \texttt{\tiny{[P1$\rangle$GM]}}
                \texttt{BOG: C, (6, 7)} \\
            }
        }
    }
    & & \\ \\

    \theutterance \stepcounter{utterance}  
    & & & \multicolumn{2}{p{0.3\linewidth}}{
        \cellcolor[rgb]{0.9,0.9,0.9}{
            \makecell[{{p{\linewidth}}}]{
                \texttt{\tiny{[GM$|$GM]}}
                \texttt{BOG: C, (6, 7)} \\
            }
        }
    }
    & & \\ \\

    \theutterance \stepcounter{utterance}  
    & & & \multicolumn{2}{p{0.3\linewidth}}{
        \cellcolor[rgb]{0.9,0.9,0.9}{
            \makecell[{{p{\linewidth}}}]{
                \texttt{\tiny{[GM$|$GM]}}
                \texttt{Pionós: níl (6,7) folamh, ach tá '│' ann.} \\
\\ 
\texttt{Tá 4 as 12 pionós carntha agaibh le chéile.} \\
\texttt{Déan iarracht arís, le do thoil!} \\
            }
        }
    }
    & & \\ \\

    \theutterance \stepcounter{utterance}  
    & & \multicolumn{4}{p{0.6\linewidth}}{
        \cellcolor[rgb]{0.9,0.9,0.9}{
            \makecell[{{p{\linewidth}}}]{
                \texttt{\tiny{[P1$\langle$GM]}}
                \texttt{Pionós: níl (6,7) folamh, ach tá '│' ann.} \\
\\ 
\texttt{Tá 4 as 12 pionós carntha agaibh le chéile.} \\
\texttt{Déan iarracht arís, le do thoil!} \\
            }
        }
    }
    & \\ \\

    \theutterance \stepcounter{utterance}  
    & \multicolumn{4}{p{0.6\linewidth}}{
        \cellcolor[rgb]{0.9,0.9,0.9}{
            \makecell[{{p{\linewidth}}}]{
                \texttt{\tiny{[P1$\rangle$GM]}}
                \texttt{ABAIR: Bhí fráma ag (6,7) ar mo chlár, faraor. Níl an sprioc a shocraíomar dlisteanach. Tá beagán ama fágach, mar sin caithfimid réiteach gasta a fháil. Seo na spotaí folamh glan atá fágtha agam anois: (1,2), (2,2), (4,2), (6,2), (7,2), (1,6), (4,6). An féidir leat bogadh anseo?} \\
            }
        }
    }
    & & \\ \\

    \theutterance \stepcounter{utterance}  
    & & & \multicolumn{2}{p{0.3\linewidth}}{
        \cellcolor[rgb]{0.9,0.9,0.9}{
            \makecell[{{p{\linewidth}}}]{
                \texttt{\tiny{[GM$|$GM]}}
                \texttt{ABAIR: Bhí fráma ag (6,7) ar mo chlár, faraor. Níl an sprioc a shocraíomar dlisteanach. Tá beagán ama fágach, mar sin caithfimid réiteach gasta a fháil. Seo na spotaí folamh glan atá fágtha agam anois: (1,2), (2,2), (4,2), (6,2), (7,2), (1,6), (4,6). An féidir leat bogadh anseo?} \\
            }
        }
    }
    & & \\ \\

    \theutterance \stepcounter{utterance}  
    & & \multicolumn{4}{p{0.6\linewidth}}{
        \cellcolor[rgb]{0.9,0.9,0.9}{
            \makecell[{{p{\linewidth}}}]{
                \texttt{\tiny{[GM$\rangle$P2]}}
                \texttt{Bogadh 'H' go (2,1) go rathúil. Seo mar atá do ghreille nuashonraithe:} \\
\texttt{```} \\
\texttt{ 1234567} \\
\texttt{╔═══════╗} \\
\texttt{║◌H◌◌T◌◌║ 1} \\
\texttt{║◌◌I◌◌◌◌║ 2} \\
\texttt{║◌◌◌◌◌┌─╢ 3} \\
\texttt{╟──┐◌◌│◌║ 4} \\
\texttt{║◌◌├──┼─╢ 5} \\
\texttt{║◌W│◌◌│C║ 6} \\
\texttt{║◌◌│◌◌│◌║ 7} \\
\texttt{╚══╧══╧═╝} \\
\texttt{```} \\
\texttt{Tá sibh ag imirt babhta 12 faoi láthair as uasmhéid 20.} \\
\texttt{Tá 4 as 12 pionós carntha agaibh le chéile.} \\
\texttt{Sheol an t{-}imreoir eile an teachtaireacht seo:} \\
\\ 
\\ 
\texttt{"Bhí fráma ag (6,7) ar mo chlár, faraor. Níl an sprioc a shocraíomar dlisteanach. Tá beagán ama fágach, mar sin caithfimid réiteach gasta a fháil. Seo na spotaí folamh glan atá fágtha agam anois: (1,2), (2,2), (4,2), (6,2), (7,2), (1,6), (4,6). An féidir leat bogadh anseo?"} \\
\\ 
\\ 
\texttt{Cad é do chéad ordú eile?} \\
            }
        }
    }
    & \\ \\

    \theutterance \stepcounter{utterance}  
    & & & \multicolumn{4}{p{0.6\linewidth}}{
        \cellcolor[rgb]{0.9,0.9,0.9}{
            \makecell[{{p{\linewidth}}}]{
                \texttt{\tiny{[GM$\langle$P2]}}
                \texttt{ABAIR: D'éirigh liom H a bhogadh go (2,1). Ó fhéachaim ar do chuid spotaí folamh, is éard atá fágtha agam le bogadh ná C ó (7,6) go (6,7) — ach tá fráma ansin agamsa freisin. Mar sin is í (7,7) an t{-}aon rogha eile atá fágtha. An bhfuil (7,7) folamh glan agatsa?} \\
            }
        }
    }
     \\ \\

    \theutterance \stepcounter{utterance}  
    & & & \multicolumn{2}{p{0.3\linewidth}}{
        \cellcolor[rgb]{0.9,0.9,0.9}{
            \makecell[{{p{\linewidth}}}]{
                \texttt{\tiny{[GM$|$GM]}}
                \texttt{ABAIR: D'éirigh liom H a bhogadh go (2,1). Ó fhéachaim ar do chuid spotaí folamh, is éard atá fágtha agam le bogadh ná C ó (7,6) go (6,7) — ach tá fráma ansin agamsa freisin. Mar sin is í (7,7) an t{-}aon rogha eile atá fágtha. An bhfuil (7,7) folamh glan agatsa?} \\
            }
        }
    }
    & & \\ \\

    \theutterance \stepcounter{utterance}  
    & & \multicolumn{4}{p{0.6\linewidth}}{
        \cellcolor[rgb]{0.9,0.9,0.9}{
            \makecell[{{p{\linewidth}}}]{
                \texttt{\tiny{[P1$\langle$GM]}}
                \texttt{Cuireadh do theachtaireacht ar aghaidh chuig an imreoir eile.} \\
\\ 
\\ 
\texttt{Tá sibh ag imirt babhta 12 faoi láthair as uasmhéid 20.} \\
\texttt{Tá 4 as 12 pionós carntha agaibh le chéile.} \\
\texttt{Sheol an t{-}imreoir eile an teachtaireacht seo:} \\
\\ 
\\ 
\texttt{"D'éirigh liom H a bhogadh go (2,1). Ó fhéachaim ar do chuid spotaí folamh, is éard atá fágtha agam le bogadh ná C ó (7,6) go (6,7) — ach tá fráma ansin agamsa freisin. Mar sin is í (7,7) an t{-}aon rogha eile atá fágtha. An bhfuil (7,7) folamh glan agatsa?"} \\
\\ 
\\ 
\texttt{Cad é do chéad ordú eile?} \\
            }
        }
    }
    & \\ \\

    \theutterance \stepcounter{utterance}  
    & \multicolumn{4}{p{0.6\linewidth}}{
        \cellcolor[rgb]{0.9,0.9,0.9}{
            \makecell[{{p{\linewidth}}}]{
                \texttt{\tiny{[P1$\rangle$GM]}}
                \texttt{ABAIR: Ní folamh glan atá (7,7) agamsa, tá ball den fhráma '╝' ann. Tá an cluiche ag teacht chun deiridh agus tá frámaí ag cur isteach orainn. Tá 4 cinn de mo spotaí folamh fágtha: (1,2), (2,2), (6,2), (7,2). An bhfuil aon cheann de seo folamh glan agatsa? Caithfimid C a chur in áit atá folamh ar an dá chlár.} \\
            }
        }
    }
    & & \\ \\

    \theutterance \stepcounter{utterance}  
    & & & \multicolumn{2}{p{0.3\linewidth}}{
        \cellcolor[rgb]{0.9,0.9,0.9}{
            \makecell[{{p{\linewidth}}}]{
                \texttt{\tiny{[GM$|$GM]}}
                \texttt{ABAIR: Ní folamh glan atá (7,7) agamsa, tá ball den fhráma '╝' ann. Tá an cluiche ag teacht chun deiridh agus tá frámaí ag cur isteach orainn. Tá 4 cinn de mo spotaí folamh fágtha: (1,2), (2,2), (6,2), (7,2). An bhfuil aon cheann de seo folamh glan agatsa? Caithfimid C a chur in áit atá folamh ar an dá chlár.} \\
            }
        }
    }
    & & \\ \\

    \theutterance \stepcounter{utterance}  
    & & \multicolumn{4}{p{0.6\linewidth}}{
        \cellcolor[rgb]{0.9,0.9,0.9}{
            \makecell[{{p{\linewidth}}}]{
                \texttt{\tiny{[GM$\rangle$P2]}}
                \texttt{Cuireadh do theachtaireacht ar aghaidh chuig an imreoir eile.} \\
\\ 
\\ 
\texttt{Tá sibh ag imirt babhta 13 faoi láthair as uasmhéid 20.} \\
\texttt{Tá 4 as 12 pionós carntha agaibh le chéile.} \\
\texttt{Sheol an t{-}imreoir eile an teachtaireacht seo:} \\
\\ 
\\ 
\texttt{"Ní folamh glan atá (7,7) agamsa, tá ball den fhráma '╝' ann. Tá an cluiche ag teacht chun deiridh agus tá frámaí ag cur isteach orainn. Tá 4 cinn de mo spotaí folamh fágtha: (1,2), (2,2), (6,2), (7,2). An bhfuil aon cheann de seo folamh glan agatsa? Caithfimid C a chur in áit atá folamh ar an dá chlár."} \\
\\ 
\\ 
\texttt{Cad é do chéad ordú eile?} \\
            }
        }
    }
    & \\ \\

    \theutterance \stepcounter{utterance}  
    & & & \multicolumn{4}{p{0.6\linewidth}}{
        \cellcolor[rgb]{0.9,0.9,0.9}{
            \makecell[{{p{\linewidth}}}]{
                \texttt{\tiny{[GM$\langle$P2]}}
                \texttt{ABAIR: Tá (6,2) agus (7,2) beirt folamh glan agamsa! Mholaim C a chur ag (6,2). An bhfuil (7,2) níos fearr duitse?} \\
            }
        }
    }
     \\ \\

    \theutterance \stepcounter{utterance}  
    & & & \multicolumn{2}{p{0.3\linewidth}}{
        \cellcolor[rgb]{0.9,0.9,0.9}{
            \makecell[{{p{\linewidth}}}]{
                \texttt{\tiny{[GM$|$GM]}}
                \texttt{ABAIR: Tá (6,2) agus (7,2) beirt folamh glan agamsa! Mholaim C a chur ag (6,2). An bhfuil (7,2) níos fearr duitse?} \\
            }
        }
    }
    & & \\ \\

    \theutterance \stepcounter{utterance}  
    & & \multicolumn{4}{p{0.6\linewidth}}{
        \cellcolor[rgb]{0.9,0.9,0.9}{
            \makecell[{{p{\linewidth}}}]{
                \texttt{\tiny{[P1$\langle$GM]}}
                \texttt{Cuireadh do theachtaireacht ar aghaidh chuig an imreoir eile.} \\
\\ 
\\ 
\texttt{Tá sibh ag imirt babhta 13 faoi láthair as uasmhéid 20.} \\
\texttt{Tá 4 as 12 pionós carntha agaibh le chéile.} \\
\texttt{Sheol an t{-}imreoir eile an teachtaireacht seo:} \\
\\ 
\\ 
\texttt{"Tá (6,2) agus (7,2) beirt folamh glan agamsa! Mholaim C a chur ag (6,2). An bhfuil (7,2) níos fearr duitse?"} \\
\\ 
\\ 
\texttt{Cad é do chéad ordú eile?} \\
            }
        }
    }
    & \\ \\

    \theutterance \stepcounter{utterance}  
    & \multicolumn{4}{p{0.6\linewidth}}{
        \cellcolor[rgb]{0.9,0.9,0.9}{
            \makecell[{{p{\linewidth}}}]{
                \texttt{\tiny{[P1$\rangle$GM]}}
                \texttt{BOG: C, (6, 2)} \\
            }
        }
    }
    & & \\ \\

    \theutterance \stepcounter{utterance}  
    & & & \multicolumn{2}{p{0.3\linewidth}}{
        \cellcolor[rgb]{0.9,0.9,0.9}{
            \makecell[{{p{\linewidth}}}]{
                \texttt{\tiny{[GM$|$GM]}}
                \texttt{BOG: C, (6, 2)} \\
            }
        }
    }
    & & \\ \\

    \theutterance \stepcounter{utterance}  
    & & & \multicolumn{2}{p{0.3\linewidth}}{
        \cellcolor[rgb]{0.9,0.9,0.9}{
            \makecell[{{p{\linewidth}}}]{
                \texttt{\tiny{[GM$|$GM]}}
                \texttt{C moved to (6, 2)} \\
            }
        }
    }
    & & \\ \\

    \theutterance \stepcounter{utterance}  
    & & \multicolumn{4}{p{0.6\linewidth}}{
        \cellcolor[rgb]{0.9,0.9,0.9}{
            \makecell[{{p{\linewidth}}}]{
                \texttt{\tiny{[GM$\rangle$P2]}}
                \texttt{Cuireadh do theachtaireacht ar aghaidh chuig an imreoir eile.} \\
\\ 
\\ 
\texttt{Tá sibh ag imirt babhta 14 faoi láthair as uasmhéid 20.} \\
\texttt{Tá 4 as 12 pionós carntha agaibh le chéile.} \\
\texttt{Bhog an t{-}imreoir eile réad ar a ngreille.} \\
\\ 
\\ 
\texttt{Cad é do chéad ordú eile?} \\
            }
        }
    }
    & \\ \\

    \theutterance \stepcounter{utterance}  
    & & & \multicolumn{4}{p{0.6\linewidth}}{
        \cellcolor[rgb]{0.9,0.9,0.9}{
            \makecell[{{p{\linewidth}}}]{
                \texttt{\tiny{[GM$\langle$P2]}}
                \texttt{ABAIR: Táimid chomh gar sin! Tá C le bogadh fós agam, ach ní dhéanfaidh mé é go dtí go mbeidh aontaíocht againn ar an sprioc deiridh. An bhfuil (6,2) roghnaithe agatsa, nó (7,2)? Inis dom go mbeidh deifreach!} \\
            }
        }
    }
     \\ \\

    \theutterance \stepcounter{utterance}  
    & & & \multicolumn{2}{p{0.3\linewidth}}{
        \cellcolor[rgb]{0.9,0.9,0.9}{
            \makecell[{{p{\linewidth}}}]{
                \texttt{\tiny{[GM$|$GM]}}
                \texttt{ABAIR: Táimid chomh gar sin! Tá C le bogadh fós agam, ach ní dhéanfaidh mé é go dtí go mbeidh aontaíocht againn ar an sprioc deiridh. An bhfuil (6,2) roghnaithe agatsa, nó (7,2)? Inis dom go mbeidh deifreach!} \\
            }
        }
    }
    & & \\ \\

    \theutterance \stepcounter{utterance}  
    & & \multicolumn{4}{p{0.6\linewidth}}{
        \cellcolor[rgb]{0.9,0.9,0.9}{
            \makecell[{{p{\linewidth}}}]{
                \texttt{\tiny{[P1$\langle$GM]}}
                \texttt{Bogadh 'C' go (6,2) go rathúil. Seo mar atá do ghreille nuashonraithe:} \\
\texttt{```} \\
\texttt{ 1234567} \\
\texttt{╔═══════╗} \\
\texttt{║T◌◌W◌◌◌║ 1} \\
\texttt{║◌◌◌◌◌C◌║ 2} \\
\texttt{║◌◌◌◌◌┌─╢ 3} \\
\texttt{╟──┐◌◌│H║ 4} \\
\texttt{║◌◌├──┼─╢ 5} \\
\texttt{║◌◌│I◌│◌║ 6} \\
\texttt{║◌◌│◌◌│◌║ 7} \\
\texttt{╚══╧══╧═╝} \\
\texttt{```} \\
\texttt{Tá sibh ag imirt babhta 14 faoi láthair as uasmhéid 20.} \\
\texttt{Tá 4 as 12 pionós carntha agaibh le chéile.} \\
\texttt{Sheol an t{-}imreoir eile an teachtaireacht seo:} \\
\\ 
\\ 
\texttt{"Táimid chomh gar sin! Tá C le bogadh fós agam, ach ní dhéanfaidh mé é go dtí go mbeidh aontaíocht againn ar an sprioc deiridh. An bhfuil (6,2) roghnaithe agatsa, nó (7,2)? Inis dom go mbeidh deifreach!"} \\
\\ 
\\ 
\texttt{Cad é do chéad ordú eile?} \\
            }
        }
    }
    & \\ \\

    \theutterance \stepcounter{utterance}  
    & \multicolumn{4}{p{0.6\linewidth}}{
        \cellcolor[rgb]{0.9,0.9,0.9}{
            \makecell[{{p{\linewidth}}}]{
                \texttt{\tiny{[P1$\rangle$GM]}}
                \texttt{ABAIR: D'éirigh liom C a bhogadh go (6,2) go rathúil! Tá T ag (1,1), W ag (4,1), C ag (6,2) agam anois. Ar mo chlár féin, tá I fós ag (3,6) agus H ag (7,4), ach is cuma cén áit a bhfuil siad sin agam. An bhfuil do chuid féin ailínithe?Más ailínithe, abair críochnaithe?} \\
            }
        }
    }
    & & \\ \\

    \theutterance \stepcounter{utterance}  
    & & & \multicolumn{2}{p{0.3\linewidth}}{
        \cellcolor[rgb]{0.9,0.9,0.9}{
            \makecell[{{p{\linewidth}}}]{
                \texttt{\tiny{[GM$|$GM]}}
                \texttt{ABAIR: D'éirigh liom C a bhogadh go (6,2) go rathúil! Tá T ag (1,1), W ag (4,1), C ag (6,2) agam anois. Ar mo chlár féin, tá I fós ag (3,6) agus H ag (7,4), ach is cuma cén áit a bhfuil siad sin agam. An bhfuil do chuid féin ailínithe?Más ailínithe, abair críochnaithe?} \\
            }
        }
    }
    & & \\ \\

    \theutterance \stepcounter{utterance}  
    & & \multicolumn{4}{p{0.6\linewidth}}{
        \cellcolor[rgb]{0.9,0.9,0.9}{
            \makecell[{{p{\linewidth}}}]{
                \texttt{\tiny{[GM$\rangle$P2]}}
                \texttt{Cuireadh do theachtaireacht ar aghaidh chuig an imreoir eile.} \\
\\ 
\\ 
\texttt{Tá sibh ag imirt babhta 15 faoi láthair as uasmhéid 20.} \\
\texttt{Tá 4 as 12 pionós carntha agaibh le chéile.} \\
\texttt{Sheol an t{-}imreoir eile an teachtaireacht seo:} \\
\\ 
\\ 
\texttt{"D'éirigh liom C a bhogadh go (6,2) go rathúil! Tá T ag (1,1), W ag (4,1), C ag (6,2) agam anois. Ar mo chlár féin, tá I fós ag (3,6) agus H ag (7,4), ach is cuma cén áit a bhfuil siad sin agam. An bhfuil do chuid féin ailínithe?Más ailínithe, abair críochnaithe?"} \\
\\ 
\\ 
\texttt{Cad é do chéad ordú eile?} \\
            }
        }
    }
    & \\ \\

    \theutterance \stepcounter{utterance}  
    & & & \multicolumn{4}{p{0.6\linewidth}}{
        \cellcolor[rgb]{0.9,0.9,0.9}{
            \makecell[{{p{\linewidth}}}]{
                \texttt{\tiny{[GM$\langle$P2]}}
                \texttt{BOG: C, (6, 2)} \\
            }
        }
    }
     \\ \\

    \theutterance \stepcounter{utterance}  
    & & & \multicolumn{2}{p{0.3\linewidth}}{
        \cellcolor[rgb]{0.9,0.9,0.9}{
            \makecell[{{p{\linewidth}}}]{
                \texttt{\tiny{[GM$|$GM]}}
                \texttt{BOG: C, (6, 2)} \\
            }
        }
    }
    & & \\ \\

    \theutterance \stepcounter{utterance}  
    & & & \multicolumn{2}{p{0.3\linewidth}}{
        \cellcolor[rgb]{0.9,0.9,0.9}{
            \makecell[{{p{\linewidth}}}]{
                \texttt{\tiny{[GM$|$GM]}}
                \texttt{C moved to (6, 2)} \\
            }
        }
    }
    & & \\ \\

    \theutterance \stepcounter{utterance}  
    & & \multicolumn{4}{p{0.6\linewidth}}{
        \cellcolor[rgb]{0.9,0.9,0.9}{
            \makecell[{{p{\linewidth}}}]{
                \texttt{\tiny{[P1$\langle$GM]}}
                \texttt{Cuireadh do theachtaireacht ar aghaidh chuig an imreoir eile.} \\
\\ 
\\ 
\texttt{Tá sibh ag imirt babhta 15 faoi láthair as uasmhéid 20.} \\
\texttt{Tá 4 as 12 pionós carntha agaibh le chéile.} \\
\texttt{Bhog an t{-}imreoir eile réad ar a ngreille.} \\
\\ 
\\ 
\texttt{Cad é do chéad ordú eile?} \\
            }
        }
    }
    & \\ \\

    \theutterance \stepcounter{utterance}  
    & \multicolumn{4}{p{0.6\linewidth}}{
        \cellcolor[rgb]{0.9,0.9,0.9}{
            \makecell[{{p{\linewidth}}}]{
                \texttt{\tiny{[P1$\rangle$GM]}}
                \texttt{ABAIR: Tá gach rud ailínithe agamsa. An bhfuil tú réidh?} \\
            }
        }
    }
    & & \\ \\

    \theutterance \stepcounter{utterance}  
    & & & \multicolumn{2}{p{0.3\linewidth}}{
        \cellcolor[rgb]{0.9,0.9,0.9}{
            \makecell[{{p{\linewidth}}}]{
                \texttt{\tiny{[GM$|$GM]}}
                \texttt{ABAIR: Tá gach rud ailínithe agamsa. An bhfuil tú réidh?} \\
            }
        }
    }
    & & \\ \\

    \theutterance \stepcounter{utterance}  
    & & \multicolumn{4}{p{0.6\linewidth}}{
        \cellcolor[rgb]{0.9,0.9,0.9}{
            \makecell[{{p{\linewidth}}}]{
                \texttt{\tiny{[GM$\rangle$P2]}}
                \texttt{Bogadh 'C' go (6,2) go rathúil. Seo mar atá do ghreille nuashonraithe:} \\
\texttt{```} \\
\texttt{ 1234567} \\
\texttt{╔═══════╗} \\
\texttt{║◌H◌◌T◌◌║ 1} \\
\texttt{║◌◌I◌◌C◌║ 2} \\
\texttt{║◌◌◌◌◌┌─╢ 3} \\
\texttt{╟──┐◌◌│◌║ 4} \\
\texttt{║◌◌├──┼─╢ 5} \\
\texttt{║◌W│◌◌│◌║ 6} \\
\texttt{║◌◌│◌◌│◌║ 7} \\
\texttt{╚══╧══╧═╝} \\
\texttt{```} \\
\texttt{Tá sibh ag imirt babhta 16 faoi láthair as uasmhéid 20.} \\
\texttt{Tá 4 as 12 pionós carntha agaibh le chéile.} \\
\texttt{Sheol an t{-}imreoir eile an teachtaireacht seo:} \\
\\ 
\\ 
\texttt{"Tá gach rud ailínithe agamsa. An bhfuil tú réidh?"} \\
\\ 
\\ 
\texttt{Cad é do chéad ordú eile?} \\
            }
        }
    }
    & \\ \\

    \theutterance \stepcounter{utterance}  
    & & & \multicolumn{4}{p{0.6\linewidth}}{
        \cellcolor[rgb]{0.9,0.9,0.9}{
            \makecell[{{p{\linewidth}}}]{
                \texttt{\tiny{[GM$\langle$P2]}}
                \texttt{ABAIR: críochnaithe!} \\
            }
        }
    }
     \\ \\

    \theutterance \stepcounter{utterance}  
    & & & \multicolumn{2}{p{0.3\linewidth}}{
        \cellcolor[rgb]{0.9,0.9,0.9}{
            \makecell[{{p{\linewidth}}}]{
                \texttt{\tiny{[GM$|$GM]}}
                \texttt{ABAIR: críochnaithe!} \\
            }
        }
    }
    & & \\ \\

    \theutterance \stepcounter{utterance}  
    & & & \multicolumn{2}{p{0.3\linewidth}}{
        \cellcolor[rgb]{0.9,0.9,0.9}{
            \makecell[{{p{\linewidth}}}]{
                \texttt{\tiny{[GM$|$GM]}}
                \texttt{true} \\
            }
        }
    }
    & & \\ \\

    \theutterance \stepcounter{utterance}  
    & & & \multicolumn{2}{p{0.3\linewidth}}{
        \cellcolor[rgb]{0.9,0.9,0.9}{
            \makecell[{{p{\linewidth}}}]{
                \texttt{\tiny{[GM$|$GM]}}
                \texttt{Initial states:} \\
\texttt{```} \\
\texttt{Player 1:} \\
\texttt{ 1234567} \\
\texttt{╔═══════╗} \\
\texttt{║◌◌◌◌◌◌◌║ 1} \\
\texttt{║◌◌◌◌◌◌◌║ 2} \\
\texttt{║◌TW◌◌┌─╢ 3} \\
\texttt{╟──┐◌◌│H║ 4} \\
\texttt{║C◌├──┼─╢ 5} \\
\texttt{║◌◌│I◌│◌║ 6} \\
\texttt{║◌◌│◌◌│◌║ 7} \\
\texttt{╚══╧══╧═╝} \\
\\ 
\\ 
\texttt{Player 2:} \\
\texttt{ 1234567} \\
\texttt{╔═══════╗} \\
\texttt{║◌◌◌◌TH◌║ 1} \\
\texttt{║◌◌◌◌◌◌◌║ 2} \\
\texttt{║◌◌◌◌◌┌─╢ 3} \\
\texttt{╟──┐◌◌│◌║ 4} \\
\texttt{║I◌├──┼─╢ 5} \\
\texttt{║◌W│◌◌│C║ 6} \\
\texttt{║◌◌│◌◌│◌║ 7} \\
\texttt{╚══╧══╧═╝} \\
\\ 
\texttt{```} \\
            }
        }
    }
    & & \\ \\

    \theutterance \stepcounter{utterance}  
    & & & \multicolumn{2}{p{0.3\linewidth}}{
        \cellcolor[rgb]{0.9,0.9,0.9}{
            \makecell[{{p{\linewidth}}}]{
                \texttt{\tiny{[GM$|$GM]}}
                \texttt{End states:} \\
\texttt{```} \\
\texttt{Player 1:} \\
\texttt{ 1234567} \\
\texttt{╔═══════╗} \\
\texttt{║T◌◌W◌◌◌║ 1} \\
\texttt{║◌◌◌◌◌C◌║ 2} \\
\texttt{║◌◌◌◌◌┌─╢ 3} \\
\texttt{╟──┐◌◌│H║ 4} \\
\texttt{║◌◌├──┼─╢ 5} \\
\texttt{║◌◌│I◌│◌║ 6} \\
\texttt{║◌◌│◌◌│◌║ 7} \\
\texttt{╚══╧══╧═╝} \\
\\ 
\\ 
\texttt{Player 2:} \\
\texttt{ 1234567} \\
\texttt{╔═══════╗} \\
\texttt{║◌H◌◌T◌◌║ 1} \\
\texttt{║◌◌I◌◌C◌║ 2} \\
\texttt{║◌◌◌◌◌┌─╢ 3} \\
\texttt{╟──┐◌◌│◌║ 4} \\
\texttt{║◌◌├──┼─╢ 5} \\
\texttt{║◌W│◌◌│◌║ 6} \\
\texttt{║◌◌│◌◌│◌║ 7} \\
\texttt{╚══╧══╧═╝} \\
\\ 
\texttt{```} \\
            }
        }
    }
    & & \\ \\

    \theutterance \stepcounter{utterance}  
    & & & \multicolumn{2}{p{0.3\linewidth}}{
        \cellcolor[rgb]{0.9,0.9,0.9}{
            \makecell[{{p{\linewidth}}}]{
                \texttt{\tiny{[GM$|$GM]}}
                \texttt{* success: True} \\
\texttt{* lose: False} \\
\texttt{* aborted: False} \\
\texttt{{-}{-}{-}{-}{-}{-}{-}} \\
\texttt{* Shifts: 4.00} \\
\texttt{* Max Shifts: 8.00} \\
\texttt{* Min Shifts: 4.00} \\
\texttt{* End Distance Sum: 19.34} \\
\texttt{* Init Distance Sum: 19.18} \\
\texttt{* Expected Distance Sum: 20.95} \\
\texttt{* Penalties: 4.00} \\
\texttt{* Max Penalties: 12.00} \\
\texttt{* Rounds: 16.00} \\
\texttt{* Max Rounds: 20.00} \\
\texttt{* Object Count: 5.00} \\
            }
        }
    }
    & & \\ \\

\end{supertabular}
}

\end{document}
